% Complete reviewed Unit 12; editorial master: chapters/ch12/student.md
% Markdown SHA-256: 1ac857fe82fb5d50474162487e47a794c0f2afb50230a7f3323e36b55929c54b
% One-use initial assembly followed by complete content and page review.
\newcommand{\runningtitle}{Unit 12. Tense, sequence, and voice}
% Layout inherited from the reviewed previous course; set a running title per unit.
% This is a small hand-maintained TeX layout, not a conversion pipeline.
\RequirePackage{fix-cm}
\documentclass[11pt,twoside]{article}
\usepackage{fontspec}
\setmainfont{Cambria}[
  Renderer=HarfBuzz,
  BoldFont=cambriab.ttf,
  ItalicFont=cambriai.ttf,
  BoldItalicFont=cambriaz.ttf
]
\setmonofont{lmmono10-regular.otf}[BoldFont=lmmonolt10-bold.otf]
\usepackage[paperwidth=165mm,paperheight=241mm,inner=21mm,outer=19mm,
  top=22mm,bottom=20mm,headheight=14pt,headsep=7mm,footskip=10mm]{geometry}
\usepackage{xcolor,array,booktabs,longtable,tabularx}
\usepackage{enumitem,fancyhdr}
\usepackage[unicode,hidelinks]{hyperref}
\definecolor{ink}{HTML}{000000}
\definecolor{muted}{HTML}{4C5558}
\setlength{\parindent}{0pt}
\setlength{\parskip}{3.5pt}
\AtBeginDocument{\fontsize{11}{14.3}\selectfont}
\setlength{\tabcolsep}{4pt}
\renewcommand{\arraystretch}{1.10}
\setlist[itemize]{leftmargin=1.3em,label=\textbullet,itemsep=2pt,topsep=3pt,parsep=0pt}
\setlist[enumerate]{leftmargin=1.8em,itemsep=3pt,topsep=4pt,parsep=0pt}
\widowpenalty=10000
\clubpenalty=10000
\emergencystretch=1em
\raggedbottom
\pagestyle{fancy}
\fancyhf{}
\providecommand{\runningtitle}{Old English}
\fancyhead[LE,RO]{\small\itshape \runningtitle}
\fancyhead[LO,RE]{\small Old English}
\fancyfoot[LE,RO]{\small\thepage}
\renewcommand{\headrulewidth}{0.3pt}
\newcommand{\booktitle}[2]{
  {\fontsize{9}{11}\selectfont\color{muted}#1\par}
  \vspace{3pt}
  {\fontsize{20}{24}\selectfont\color{ink}#2\par}
  \vspace{6pt}
}
\newcommand{\source}[1]{\par{\footnotesize\color{muted}#1\par}}
\newcommand{\readingtext}[1]{
  \par\addvspace{5pt}
  \begingroup\leftskip=1.5em
    \fontsize{12}{16}\selectfont #1\par
  \endgroup\addvspace{5pt}
}
\makeatletter
\renewcommand{\section}{\@startsection{section}{1}{\z@}
  {-2.4ex \@plus -0.6ex \@minus -0.2ex}{0.8ex \@plus 0.2ex}
  {\normalfont\fontsize{12}{15}\selectfont\bfseries\color{ink}}}
\renewcommand{\subsection}{\@startsection{subsection}{2}{\z@}
  {-1.8ex \@plus -0.5ex \@minus -0.2ex}{0.5ex \@plus 0.2ex}
  {\normalfont\fontsize{11}{14}\selectfont\bfseries\color{ink}}}
\makeatother

\hyphenpenalty=10000
\exhyphenpenalty=10000
\emergencystretch=3em
\newcommand{\itempara}[2]{\par\addvspace{3pt}\noindent\hangindent=1.8em\hangafter=1\makebox[1.8em][l]{#1.}#2\par}
\newcommand{\reserve}[1]{\par\penalty-100\begingroup\dimen0=\pagegoal\advance\dimen0 by-\pagetotal\ifdim\dimen0<#1\newpage\fi\endgroup}
\begin{document}
\addfontfeatures{Ligatures=TeXOff}
\raggedright
\booktitle{UNIT 12}{Tense, sequence, and voice}

% mdparagraph:P001:begin
Old English has two inflectional tenses, present and preterite. Their interpretation depends on the situation. A present may refer to a future event; a preterite may narrate an event, describe a past state, or refer to something already completed at a past viewpoint. English may consequently need a simple, perfect, pluperfect or future expression where Old English has one finite word.
% mdparagraph:P001:end

% mdparagraph:P002:begin
The twelve translated examples, exercises F, A, C, R and S, and “A house and a return” are new teaching material. Cook VI.19–21 retain all their original English prompts. The final historical reading is an exact connected extract, with separate English context, learner help and Cook’s original notes.
% mdparagraph:P002:end

% mdparagraph:P003:begin
In tables, 1, 2 and 3 denote person. Nominative, genitive, dative and accusative are written out. Repeated spellings occupy distinct grammatical positions. \textbf{Reference} means recognize the stated form or construction; it is not required in independent composition. \textbf{Supplied help} remains available while answering the associated task. Preserve written vowel lengths and the distinction between e and ę.
% mdparagraph:P003:end

\reserve{7\baselineskip}
% mdparagraph:P004:begin
\section*{1. Time in a narrative}
% mdparagraph:P004:end

% teaching: NEW-U12-GRAM-u12-narrative-time
% teaching: NEW-U12-GRAM-u12-preterite-anterior
% mdparagraph:P005:begin
A succession of preterites commonly advances the narrative. Temporal clauses, the meanings of the verbs and the surrounding account establish which event is earlier. A clause does not acquire a different inflection simply because English translates it with “had”.
% mdparagraph:P005:end

\reserve{8\baselineskip}
% example: NEW-U12-EXAMPLE01
% mdparagraph:P006:begin
Se mann cōm, and hē geseah ðæt hūs.
% mdparagraph:P006:end

% mdparagraph:P007:begin
The man came, and he saw the house.
% mdparagraph:P007:end

% mdparagraph:P008:begin
Cōm and geseah are third singular preterite indicatives. In this account the arrival precedes the seeing.
% mdparagraph:P008:end

\reserve{12\baselineskip}
% mdparagraph:P009:begin
Context: The house was finished before the man’s visit.
% mdparagraph:P009:end

\reserve{8\baselineskip}
% example: NEW-U12-EXAMPLE02
% mdparagraph:P010:begin
Se mann geseah ðæt hūs ðe his sunu geworhte.
% mdparagraph:P010:end

% mdparagraph:P011:begin
The man saw the house which his son had built.
% mdparagraph:P011:end

% mdparagraph:P012:begin
Geworhte is also a simple preterite indicative. Here the stated context places the son’s building before the man’s seeing. English “had built” makes that relation clear without changing the Old English form.
% mdparagraph:P012:end

\reserve{12\baselineskip}
% mdparagraph:P013:begin
Context: The visit is still expected.
% mdparagraph:P013:end

\reserve{8\baselineskip}
% example: NEW-U12-EXAMPLE12
% mdparagraph:P014:begin
Hē cymð.
% mdparagraph:P014:end

% mdparagraph:P015:begin
He will come.
% mdparagraph:P015:end

% mdparagraph:P016:begin
Cymð is present indicative. The context makes the coming future. Do not add willan merely because the English translation contains “will”; willan can retain its meaning “wish; intend”.
% mdparagraph:P016:end

\reserve{7\baselineskip}
% mdparagraph:P017:begin
\section*{2. Habban with a past participle}
% mdparagraph:P017:end

% teaching: NEW-U12-GRAM-u12-habban-perfect
% teaching: NEW-U12-VOC-HABBAN
% mdparagraph:P018:begin
With a transitive past participle and an accusative object, finite \textbf{habban} can present completed action and its result. The auxiliary agrees with the subject. The object is the person or thing affected by the participle. In the following construction the speaker has performed the building, and the completed house is the result.
% mdparagraph:P018:end

\reserve{8\baselineskip}
% example: NEW-U12-EXAMPLE03
% mdparagraph:P019:begin
Nū hæbbe ic ðæt hūs geworht.
% mdparagraph:P019:end

% mdparagraph:P020:begin
Now I have built the house.
% mdparagraph:P020:end

\reserve{12\baselineskip}
% mdparagraph:P021:begin
Context: The building was already complete when the son arrived.
% mdparagraph:P021:end

\reserve{8\baselineskip}
% example: NEW-U12-EXAMPLE04
% mdparagraph:P022:begin
Ðā se sunu cōm, hæfde se mann ðæt hūs geworht.
% mdparagraph:P022:end

% mdparagraph:P023:begin
When the son came, the man had built the house.
% mdparagraph:P023:end

% mdparagraph:P024:begin
Present hæbbe presents the completed work from the present viewpoint. Past hæfde places the completed work before the stated past arrival. The difference is in the finite auxiliary; geworht remains the selected participle base.
% mdparagraph:P024:end

% teaching: NEW-U12-HABBAN-PERFECT-IND
\reserve{11\baselineskip}
% mdparagraph:P025:begin
\textbf{Habban with an object and past participle: complete present and preterite indicative construction}
% mdparagraph:P025:end

% construction: NEW-U12-HABBAN-PERFECT-IND
% mdtable:T01:begin
% course-table-style: 2026-10
% course-panel:T01-P01:begin
\par\begingroup
\fontsize{10.5}{13.2}\selectfont
\setlength{\LTleft}{0pt}\setlength{\LTright}{\fill}
\noindent
\begin{tabular}{@{}>{\raggedright\arraybackslash}p{39.459mm}>{\raggedright\arraybackslash}p{39.459mm}>{\raggedright\arraybackslash}p{39.459mm}@{}}
\toprule
\textbf{Person and number} & \textbf{Present indicative} & \textbf{Preterite indicative} \\
\midrule
1 sg. & Ic hæbbe ðæt hūs geworht. & Ic hæfde ðæt hūs geworht. \\
2 sg. & Ðū hæfst ðæt hūs geworht. & Ðū hæfdest ðæt hūs geworht. \\
3 sg. & Hē hæfð ðæt hūs geworht. & Hē hæfde ðæt hūs geworht. \\
1 pl. & Wē habbað ðæt hūs geworht. & Wē hæfdon ðæt hūs geworht. \\
2 pl. & Gē habbað ðæt hūs geworht. & Gē hæfdon ðæt hūs geworht. \\
3 pl. & Hīe habbað ðæt hūs geworht. & Hīe hæfdon ðæt hūs geworht. \\
\bottomrule
\end{tabular}
\par\endgroup
% course-panel:T01-P01:end
% mdtable:T01:end
\reserve{11\baselineskip}
% mdparagraph:P026:begin
\textbf{Translations of the twelve positions}
% mdparagraph:P026:end

% mdtable:T02:begin
% course-table-style: 2026-10
% course-panel:T02-P01:begin
\par\begingroup
\fontsize{10.5}{13.2}\selectfont
\setlength{\LTleft}{0pt}\setlength{\LTright}{\fill}
\noindent
\begin{tabular}{@{}>{\raggedright\arraybackslash}p{39.459mm}>{\raggedright\arraybackslash}p{39.459mm}>{\raggedright\arraybackslash}p{39.459mm}@{}}
\toprule
\textbf{Person and number} & \textbf{Present viewpoint} & \textbf{Past viewpoint} \\
\midrule
1 sg. & I have built the house. & I had built the house. \\
2 sg. & You have built the house. & You had built the house. \\
3 sg. & He has built the house. & He had built the house. \\
1 pl. & We have built the house. & We had built the house. \\
2 pl. & You have built the house. & You had built the house. \\
3 pl. & They have built the house. & They had built the house. \\
\bottomrule
\end{tabular}
\par\endgroup
% course-panel:T02-P01:end
% mdtable:T02:end
% mdparagraph:P027:begin
This is a table of a construction in the indicative, not a new tense ending. The object ðæt hūs stays accusative even when the subject changes. The six auxiliary forms in each tense were learned earlier.
% mdparagraph:P027:end

% mdparagraph:P028:begin
Use the bare participle bases \textbf{geworht, gierwed, sętt} in the assigned compositions. Old English also has participles agreeing with the object in this construction. A properly evidenced agreeing form is a valid alternative; the bare forms here do not teach a complete participial adjective declension. Nor does this pattern establish every causative use of habban, or its use with every intransitive verb.
% mdparagraph:P028:end

\reserve{7\baselineskip}
% mdparagraph:P029:begin
\section*{3. Active clauses, passives, and other uses of be}
% mdparagraph:P029:end

% teaching: NEW-U12-GRAM-u12-active-passive
% mdparagraph:P030:begin
An active transitive clause names an agent as subject and the affected entity as object. In an agentless passive the affected entity becomes nominative subject, and the active agent is no longer named. Use the already learned passive auxiliary with the selected past participle. Do not add a new expression for a passive agent.
% mdparagraph:P030:end

\reserve{8\baselineskip}
% example: NEW-U12-EXAMPLE05
% mdparagraph:P031:begin
Se mann band ðæt hors. Ðæt hors wearð bunden.
% mdparagraph:P031:end

% mdparagraph:P032:begin
The man bound the horse. The horse was bound.
% mdparagraph:P032:end

% mdparagraph:P033:begin
The man is the active subject, and the horse is the active object. In the second clause the horse is the subject. Wearð often presents the event or entrance into a condition; the passive itself does not tell who performed the binding.
% mdparagraph:P033:end

\reserve{12\baselineskip}
% mdparagraph:P034:begin
Context: The context presents the horse in its resulting bound state; wesan does not invariably exclude an event reading.
% mdparagraph:P034:end

\reserve{8\baselineskip}
% example: NEW-U12-EXAMPLE06
% mdparagraph:P035:begin
Ðæt hors wæs bunden.
% mdparagraph:P035:end

% mdparagraph:P036:begin
The horse was bound.
% mdparagraph:P036:end

% mdparagraph:P037:begin
Here wæs presents the resulting condition. This is a contextual contrast: wesan can also occur in accounts of events, and weorðan need not always be translated literally “became”.
% mdparagraph:P037:end

% teaching: NEW-U12-GRAM-u12-wesan-intransitive-perfect
% teaching: NEW-U12-VOC-WESAN
% mdparagraph:P038:begin
For recognition, \textbf{wesan + gewiten}, the past participle of gewītan, can express a completed departure. The subject is the person who departs, not a patient acted on by an unnamed agent.
% mdparagraph:P038:end

\reserve{12\baselineskip}
% mdparagraph:P039:begin
Context: In a past account, the man’s departure precedes the next event.
% mdparagraph:P039:end

\reserve{8\baselineskip}
% example: NEW-U12-EXAMPLE07
% mdparagraph:P040:begin
Se mann wæs gewiten.
% mdparagraph:P040:end

% mdparagraph:P041:begin
The man had departed.
% mdparagraph:P041:end

% teaching: NEW-U12-GRAM-u12-wesan-present-participle
% mdparagraph:P042:begin
Also recognize \textbf{wesan + sittende}, the present participle of sittan. The context can support a continuing activity or state. English may use a simple past or a progressive; the Old English form does not impose a modern English progressive category.
% mdparagraph:P042:end

\reserve{12\baselineskip}
% mdparagraph:P043:begin
Context: The account describes a period during which the man was seated.
% mdparagraph:P043:end

\reserve{8\baselineskip}
% example: NEW-U12-EXAMPLE08
% mdparagraph:P044:begin
Se mann wæs sittende.
% mdparagraph:P044:end

% mdparagraph:P045:begin
The man sat /\allowbreak{} was sitting.
% mdparagraph:P045:end

% mdparagraph:P046:begin
These two uses are for recognition in the assigned examples and reading. Do not construct an unrestricted intransitive perfect or progressive system from them.
% mdparagraph:P046:end

\reserve{7\baselineskip}
% mdparagraph:P047:begin
\section*{4. Time in reported speech}
% mdparagraph:P047:end

% teaching: NEW-U12-GRAM-u12-reported-sequence
% mdparagraph:P048:begin
A reported statement has its own temporal viewpoint. Ask whether the reported event is simultaneous with the saying, already completed before it, or intended afterward. The earlier reporting subjunctive remains available. A subjunctive does not mean that the report is false, and a past reporting verb does not require mechanical replacement of every dependent tense.
% mdparagraph:P048:end

\reserve{12\baselineskip}
% mdparagraph:P049:begin
Context: In each report, the presence is simultaneous with the saying.
% mdparagraph:P049:end

\reserve{8\baselineskip}
% example: NEW-U12-EXAMPLE09
% mdparagraph:P050:begin
Hē sægð ðæt hē sīe hēr. Hē sægde ðæt hē wǣre hēr.
% mdparagraph:P050:end

% mdparagraph:P051:begin
He says that he is here. He said that he was here.
% mdparagraph:P051:end

% mdparagraph:P052:begin
Sīe is present subjunctive and wǣre preterite subjunctive. Each places the reported presence at the time of its own saying. The past report does not by itself say that the person is absent now.
% mdparagraph:P052:end

\reserve{12\baselineskip}
% mdparagraph:P053:begin
Context: The building precedes the past saying.
% mdparagraph:P053:end

\reserve{8\baselineskip}
% example: NEW-U12-EXAMPLE10
% mdparagraph:P054:begin
Hē sægde ðæt hē ðæt hūs geworht hæfde.
% mdparagraph:P054:end

% mdparagraph:P055:begin
He said that he had built the house.
% mdparagraph:P055:end

% mdparagraph:P056:begin
In this assigned report, hæfde is preterite subjunctive singular. It has the same spelling as first and third singular indicative hæfde. The reported-statement construction determines the mood; the spelling alone does not. Geworht with hæfde places the building before the past saying. The complete table in §2 remains labelled indicative.
% mdparagraph:P056:end

\reserve{12\baselineskip}
% mdparagraph:P057:begin
Context: The speaker reports an intention; the sentence does not establish whether the arrival occurred.
% mdparagraph:P057:end

\reserve{8\baselineskip}
% example: NEW-U12-EXAMPLE11
% mdparagraph:P058:begin
Hē sægde ðæt hē wolde cuman.
% mdparagraph:P058:end

% mdparagraph:P059:begin
He said that he intended to come.
% mdparagraph:P059:end

% mdparagraph:P060:begin
Wolde with a plain infinitive reports a wish or intention viewed from the past. The intended coming is later than that intention. The sentence does not assert that the coming took place. This is the use of willan required here, not a general future auxiliary.
% mdparagraph:P060:end

\reserve{7\baselineskip}
% mdparagraph:P061:begin
\section*{5. Verb systems and their uses}
% mdparagraph:P061:end

% mdparagraph:P062:begin
The following indicative tables include all three persons in both numbers for each of the two tenses. Their lexical uses are stated beside them. Forms of an additional mood or a nonfinite construction are learned separately, even when a spelling happens to coincide.
% mdparagraph:P062:end

% mdparagraph:P063:begin
In the ordinary strong pattern, present singular endings are -e, -st, -ð and the plural is -að, with the stated stem changes and contractions. Preterite first and third singular share one stem; second singular uses the plural stem plus -e, and all plural persons use -on. Learn the contracted flēon distribution directly.
% mdparagraph:P063:end

% teaching: NEW-U12-GRAM-u12-adrifan-indicative
% teaching: NEW-U12-VOC-ADRIFAN
% mdparagraph:P064:begin
\textbf{ādrīfan: drive out; expel.} Ādrīfan takes an accusative person or group expelled. The earlier fram + dative can identify the place from which they are driven. Present ī contrasts with singular past ā and the short i of second singular and plural past forms.
% mdparagraph:P064:end

% teaching: NEW-U12-PAR-adrifan-IND
\reserve{11\baselineskip}
% mdparagraph:P065:begin
\textbf{ādrīfan; strong verb: present and preterite indicative}
% mdparagraph:P065:end

% table: NEW-U12-PAR-adrifan-IND
% mdtable:T03:begin
% course-table-style: 2026-10
% course-panel:T03-P01:begin
\par\begingroup
\fontsize{10.5}{13.2}\selectfont
\setlength{\LTleft}{0pt}\setlength{\LTright}{\fill}
\noindent
\begin{tabular}{@{}>{\raggedright\arraybackslash}p{39.459mm}>{\raggedright\arraybackslash}p{39.459mm}>{\raggedright\arraybackslash}p{39.459mm}@{}}
\toprule
\textbf{Person and number} & \textbf{Present} & \textbf{Preterite} \\
\midrule
1 sg. & ādrīfe & ādrāf \\
2 sg. & ādrīfst & ādrife \\
3 sg. & ādrīfð & ādrāf \\
1–3 pl. & ādrīfað & ādrifon \\
\bottomrule
\end{tabular}
\par\endgroup
% course-panel:T03-P01:end
% mdtable:T03:end
% teaching: NEW-U12-GRAM-u12-gehatan-indicative
% teaching: NEW-U12-VOC-GEHATAN
% mdparagraph:P066:begin
\textbf{gehātan: promise.} Gehātan takes the accusative thing promised and, when stated, a dative recipient. Learn hā- in the present beside hǣt- in second/\allowbreak{}third singular and hēt- in the preterite. The prefix ge- remains. No additional calling or commanding sense is assigned.
% mdparagraph:P066:end

% teaching: NEW-U12-PAR-gehatan-IND
\reserve{11\baselineskip}
% mdparagraph:P067:begin
\textbf{gehātan; strong verb: present and preterite indicative}
% mdparagraph:P067:end

% table: NEW-U12-PAR-gehatan-IND
% mdtable:T04:begin
% course-table-style: 2026-10
% course-panel:T04-P01:begin
\par\begingroup
\fontsize{10.5}{13.2}\selectfont
\setlength{\LTleft}{0pt}\setlength{\LTright}{\fill}
\noindent
\begin{tabular}{@{}>{\raggedright\arraybackslash}p{39.459mm}>{\raggedright\arraybackslash}p{39.459mm}>{\raggedright\arraybackslash}p{39.459mm}@{}}
\toprule
\textbf{Person and number} & \textbf{Present} & \textbf{Preterite} \\
\midrule
1 sg. & gehāte & gehēt \\
2 sg. & gehǣtst & gehēte \\
3 sg. & gehǣtt & gehēt \\
1–3 pl. & gehātað & gehēton \\
\bottomrule
\end{tabular}
\par\endgroup
% course-panel:T04-P01:end
% mdtable:T04:end
% teaching: NEW-U12-GRAM-u12-fleon-indicative
% teaching: NEW-U12-VOC-FLEON
% mdparagraph:P068:begin
\textbf{flēon: flee; flee from; avoid.} Flēon can stand without an object or take the accusative person or thing avoided. It is distinct from flēogan, fly. Keep flēo-, flīeh-, flēah- and flug- in their printed positions; do not form its contracted paradigm from the infinitive by adding ordinary endings.
% mdparagraph:P068:end

% teaching: NEW-U12-PAR-fleon-IND
\reserve{11\baselineskip}
% mdparagraph:P069:begin
\textbf{flēon; contracted strong verb: present and preterite indicative}
% mdparagraph:P069:end

% table: NEW-U12-PAR-fleon-IND
% mdtable:T05:begin
% course-table-style: 2026-10
% course-panel:T05-P01:begin
\par\begingroup
\fontsize{10.5}{13.2}\selectfont
\setlength{\LTleft}{0pt}\setlength{\LTright}{\fill}
\noindent
\begin{tabular}{@{}>{\raggedright\arraybackslash}p{39.459mm}>{\raggedright\arraybackslash}p{39.459mm}>{\raggedright\arraybackslash}p{39.459mm}@{}}
\toprule
\textbf{Person and number} & \textbf{Present} & \textbf{Preterite} \\
\midrule
1 sg. & flēo & flēah \\
2 sg. & flīehst & fluge \\
3 sg. & flīehð & flēah \\
1–3 pl. & flēoð & flugon \\
\bottomrule
\end{tabular}
\par\endgroup
% course-panel:T05-P01:end
% mdtable:T05:end
% teaching: NEW-U12-GRAM-u12-gemetan-indicative
% teaching: NEW-U12-VOC-GEMETAN
% mdparagraph:P070:begin
Weak I gemētan and grētan retain ē. Their t meets the third-singular dental in -tt, and their preterite has -tte, -ttest, -tton. Ēhtan and cȳðan have their own stated dental combinations.
% mdparagraph:P070:end

% mdparagraph:P071:begin
\textbf{gemētan: find; meet.} Gemētan takes the accusative person or thing found. Distinguish preterite singular gemētte from third singular present gemētt.
% mdparagraph:P071:end

% teaching: NEW-U12-PAR-gemetan-IND
\reserve{11\baselineskip}
% mdparagraph:P072:begin
\textbf{gemētan; weak verb: present and preterite indicative}
% mdparagraph:P072:end

% table: NEW-U12-PAR-gemetan-IND
% mdtable:T06:begin
% course-table-style: 2026-10
% course-panel:T06-P01:begin
\par\begingroup
\fontsize{10.5}{13.2}\selectfont
\setlength{\LTleft}{0pt}\setlength{\LTright}{\fill}
\noindent
\begin{tabular}{@{}>{\raggedright\arraybackslash}p{39.459mm}>{\raggedright\arraybackslash}p{39.459mm}>{\raggedright\arraybackslash}p{39.459mm}@{}}
\toprule
\textbf{Person and number} & \textbf{Present} & \textbf{Preterite} \\
\midrule
1 sg. & gemēte & gemētte \\
2 sg. & gemētst & gemēttest \\
3 sg. & gemētt & gemētte \\
1–3 pl. & gemētað & gemētton \\
\bottomrule
\end{tabular}
\par\endgroup
% course-panel:T06-P01:end
% mdtable:T06:end
% teaching: NEW-U12-GRAM-u12-geopenian-indicative
% teaching: NEW-U12-VOC-GEOPENIAN
% mdparagraph:P073:begin
Weak II geopenian and leornian have present -ie, -ast, -að, -iað and preterite -ode, -odest, -ode, -odon. The plural spelling is shared by all persons.
% mdparagraph:P073:end

% mdparagraph:P074:begin
\textbf{geopenian: open.} Geopenian takes the accusative thing opened, including a mouth or eyes. The table establishes this transitive use.
% mdparagraph:P074:end

% teaching: NEW-U12-PAR-geopenian-IND
\reserve{11\baselineskip}
% mdparagraph:P075:begin
\textbf{geopenian; weak verb: present and preterite indicative}
% mdparagraph:P075:end

% table: NEW-U12-PAR-geopenian-IND
% mdtable:T07:begin
% course-table-style: 2026-10
% course-panel:T07-P01:begin
\par\begingroup
\fontsize{10.5}{13.2}\selectfont
\setlength{\LTleft}{0pt}\setlength{\LTright}{\fill}
\noindent
\begin{tabular}{@{}>{\raggedright\arraybackslash}p{39.459mm}>{\raggedright\arraybackslash}p{39.459mm}>{\raggedright\arraybackslash}p{39.459mm}@{}}
\toprule
\textbf{Person and number} & \textbf{Present} & \textbf{Preterite} \\
\midrule
1 sg. & geopenie & geopenode \\
2 sg. & geopenast & geopenodest \\
3 sg. & geopenað & geopenode \\
1–3 pl. & geopeniað & geopenodon \\
\bottomrule
\end{tabular}
\par\endgroup
% course-panel:T07-P01:end
% mdtable:T07:end
% teaching: NEW-U12-GRAM-u12-gretan-indicative
% teaching: NEW-U12-VOC-GRETAN
% mdparagraph:P076:begin
\textbf{grētan: greet; take leave of.} Grētan takes an accusative person greeted or taken leave of. It also takes the native land as object in the assigned leave-taking expression. This weak verb is distinct from the strong homonym meaning weep.
% mdparagraph:P076:end

% teaching: NEW-U12-PAR-gretan-IND
\reserve{11\baselineskip}
% mdparagraph:P077:begin
\textbf{grētan; weak verb: present and preterite indicative}
% mdparagraph:P077:end

% table: NEW-U12-PAR-gretan-IND
% mdtable:T08:begin
% course-table-style: 2026-10
% course-panel:T08-P01:begin
\par\begingroup
\fontsize{10.5}{13.2}\selectfont
\setlength{\LTleft}{0pt}\setlength{\LTright}{\fill}
\noindent
\begin{tabular}{@{}>{\raggedright\arraybackslash}p{39.459mm}>{\raggedright\arraybackslash}p{39.459mm}>{\raggedright\arraybackslash}p{39.459mm}@{}}
\toprule
\textbf{Person and number} & \textbf{Present} & \textbf{Preterite} \\
\midrule
1 sg. & grēte & grētte \\
2 sg. & grētst & grēttest \\
3 sg. & grētt & grētte \\
1–3 pl. & grētað & grētton \\
\bottomrule
\end{tabular}
\par\endgroup
% course-panel:T08-P01:end
% mdtable:T08:end
% teaching: NEW-U12-GRAM-u12-cythan-indicative
% teaching: NEW-U12-VOC-CYTHAN
% mdparagraph:P078:begin
\textbf{cȳðan: announce; make known.} Cȳðan takes an accusative thing announced or a ðæt content clause; a recipient, when expressed, is dative. The selected present forms are cȳst and cȳð; learn their dental combinations directly. Analogical cȳðst follows the documented retained-dental pattern in the second singular. The selected preterite has cȳðd-; cȳdd- is an attested alternative.
% mdparagraph:P078:end

% teaching: NEW-U12-PAR-cythan-IND
\reserve{11\baselineskip}
% mdparagraph:P079:begin
\textbf{cȳðan; weak verb: present and preterite indicative}
% mdparagraph:P079:end

% table: NEW-U12-PAR-cythan-IND
% mdtable:T09:begin
% course-table-style: 2026-10
% course-panel:T09-P01:begin
\par\begingroup
\fontsize{10.5}{13.2}\selectfont
\setlength{\LTleft}{0pt}\setlength{\LTright}{\fill}
\noindent
\begin{tabular}{@{}>{\raggedright\arraybackslash}p{39.459mm}>{\raggedright\arraybackslash}p{39.459mm}>{\raggedright\arraybackslash}p{39.459mm}@{}}
\toprule
\textbf{Person and number} & \textbf{Present} & \textbf{Preterite} \\
\midrule
1 sg. & cȳðe & cȳðde \\
2 sg. & cȳst & cȳðdest \\
3 sg. & cȳð & cȳðde \\
1–3 pl. & cȳðað & cȳðdon \\
\bottomrule
\end{tabular}
\par\endgroup
% course-panel:T09-P01:end
% mdtable:T09:end
% teaching: NEW-U12-GRAM-u12-ehtan-indicative
% teaching: NEW-U12-VOC-EHTAN
% mdparagraph:P080:begin
\textbf{ēhtan: pursue; persecute.} Use the accusative person pursued in the assigned construction, directly supported by the sheep-and-wolves Gospel passage. Genitive government also occurs and is valid with supported forms. Preserve ēht in third singular present, and ēhte in the preterite.
% mdparagraph:P080:end

% teaching: NEW-U12-PAR-ehtan-IND
\reserve{11\baselineskip}
% mdparagraph:P081:begin
\textbf{ēhtan; weak verb: present and preterite indicative}
% mdparagraph:P081:end

% table: NEW-U12-PAR-ehtan-IND
% mdtable:T10:begin
% course-table-style: 2026-10
% course-panel:T10-P01:begin
\par\begingroup
\fontsize{10.5}{13.2}\selectfont
\setlength{\LTleft}{0pt}\setlength{\LTright}{\fill}
\noindent
\begin{tabular}{@{}>{\raggedright\arraybackslash}p{39.459mm}>{\raggedright\arraybackslash}p{39.459mm}>{\raggedright\arraybackslash}p{39.459mm}@{}}
\toprule
\textbf{Person and number} & \textbf{Present} & \textbf{Preterite} \\
\midrule
1 sg. & ēhte & ēhte \\
2 sg. & ēhtst & ēhtest \\
3 sg. & ēht & ēhte \\
1–3 pl. & ēhtað & ēhton \\
\bottomrule
\end{tabular}
\par\endgroup
% course-panel:T10-P01:end
% mdtable:T10:end
% teaching: NEW-U12-GRAM-u12-agan-indicative
% teaching: NEW-U12-VOC-AGAN
% mdparagraph:P082:begin
\textbf{āgan: own; possess.} Āgan takes an accusative possession. Its preterite-present pattern has āh/\allowbreak{}āhst/\allowbreak{}āh and plural āgon in the present, āhte/\allowbreak{}āhtest/\allowbreak{}āhte and plural āhton in the preterite. This is the positive indicative only. The supplied negative nāh below does not teach a complete negative system.
% mdparagraph:P082:end

% teaching: NEW-U12-PAR-agan-IND
\reserve{11\baselineskip}
% mdparagraph:P083:begin
\textbf{āgan; preterite-present verb: present and preterite indicative}
% mdparagraph:P083:end

% table: NEW-U12-PAR-agan-IND
% mdtable:T11:begin
% course-table-style: 2026-10
% course-panel:T11-P01:begin
\par\begingroup
\fontsize{10.5}{13.2}\selectfont
\setlength{\LTleft}{0pt}\setlength{\LTright}{\fill}
\noindent
\begin{tabular}{@{}>{\raggedright\arraybackslash}p{39.459mm}>{\raggedright\arraybackslash}p{39.459mm}>{\raggedright\arraybackslash}p{39.459mm}@{}}
\toprule
\textbf{Person and number} & \textbf{Present} & \textbf{Preterite} \\
\midrule
1 sg. & āh & āhte \\
2 sg. & āhst & āhtest \\
3 sg. & āh & āhte \\
1–3 pl. & āgon & āhton \\
\bottomrule
\end{tabular}
\par\endgroup
% course-panel:T11-P01:end
% mdtable:T11:end
% teaching: NEW-U12-GRAM-u12-leornian-indicative
% teaching: NEW-U12-VOC-LEORNIAN
% mdparagraph:P084:begin
\textbf{leornian: learn.} An expressed thing learned is accusative. The person from whom one learns is introduced by æt + dative. The phrase learn from me can stand without an expressed object.
% mdparagraph:P084:end

% teaching: NEW-U12-PAR-leornian-IND
\reserve{11\baselineskip}
% mdparagraph:P085:begin
\textbf{leornian; weak verb: present and preterite indicative}
% mdparagraph:P085:end

% table: NEW-U12-PAR-leornian-IND
% mdtable:T12:begin
% course-table-style: 2026-10
% course-panel:T12-P01:begin
\par\begingroup
\fontsize{10.5}{13.2}\selectfont
\setlength{\LTleft}{0pt}\setlength{\LTright}{\fill}
\noindent
\begin{tabular}{@{}>{\raggedright\arraybackslash}p{39.459mm}>{\raggedright\arraybackslash}p{39.459mm}>{\raggedright\arraybackslash}p{39.459mm}@{}}
\toprule
\textbf{Person and number} & \textbf{Present} & \textbf{Preterite} \\
\midrule
1 sg. & leornie & leornode \\
2 sg. & leornast & leornodest \\
3 sg. & leornað & leornode \\
1–3 pl. & leorniað & leornodon \\
\bottomrule
\end{tabular}
\par\endgroup
% course-panel:T12-P01:end
% mdtable:T12:end
\reserve{7\baselineskip}
% mdparagraph:P086:begin
\subsection*{Commands}
% mdparagraph:P086:end

% mdparagraph:P087:begin
A singular command addresses one person; a plural command addresses two or more. Both members of each new pair are given. The selected strong long stems and weak I long stems have a bare singular; flēon has flēoh. Weak II uses -a in the singular and -iað in the plural. These commands add no new subjunctive paradigms.
% mdparagraph:P087:end

% teaching: NEW-U12-GRAM-u12-sendan-imperative
% teaching: NEW-U12-PAR-sendan-IMP
\reserve{7\baselineskip}
% mdparagraph:P088:begin
\textbf{sęndan; weak verb: singular and plural commands}
% mdparagraph:P088:end

% table: NEW-U12-PAR-sendan-IMP
% mdtable:T13:begin
% course-table-style: 2026-10
% course-panel:T13-P01:begin
\par\begingroup
\fontsize{10.5}{13.2}\selectfont
\setlength{\LTleft}{0pt}\setlength{\LTright}{\fill}
\noindent
\begin{tabular}{@{}>{\raggedright\arraybackslash}p{63.342mm}>{\raggedright\arraybackslash}p{57.846mm}@{}}
\toprule
\textbf{Person and number} & \textbf{Command} \\
\midrule
2 sg. & sęnd \\
2 pl. & sęndað \\
\bottomrule
\end{tabular}
\par\endgroup
% course-panel:T13-P01:end
% mdtable:T13:end
% teaching: NEW-U12-GRAM-u12-bringan-imperative
% teaching: NEW-U12-PAR-bringan-IMP
\reserve{7\baselineskip}
% mdparagraph:P089:begin
\textbf{bringan; mixed verb: singular and plural commands}
% mdparagraph:P089:end

% table: NEW-U12-PAR-bringan-IMP
% mdtable:T14:begin
% course-table-style: 2026-10
% course-panel:T14-P01:begin
\par\begingroup
\fontsize{10.5}{13.2}\selectfont
\setlength{\LTleft}{0pt}\setlength{\LTright}{\fill}
\noindent
\begin{tabular}{@{}>{\raggedright\arraybackslash}p{63.342mm}>{\raggedright\arraybackslash}p{57.846mm}@{}}
\toprule
\textbf{Person and number} & \textbf{Command} \\
\midrule
2 sg. & bring \\
2 pl. & bringað \\
\bottomrule
\end{tabular}
\par\endgroup
% course-panel:T14-P01:end
% mdtable:T14:end
% teaching: NEW-U12-GRAM-u12-adrifan-imperative
% teaching: NEW-U12-PAR-adrifan-IMP
\reserve{7\baselineskip}
% mdparagraph:P090:begin
\textbf{ādrīfan; strong verb: singular and plural commands}
% mdparagraph:P090:end

% table: NEW-U12-PAR-adrifan-IMP
% mdtable:T15:begin
% course-table-style: 2026-10
% course-panel:T15-P01:begin
\par\begingroup
\fontsize{10.5}{13.2}\selectfont
\setlength{\LTleft}{0pt}\setlength{\LTright}{\fill}
\noindent
\begin{tabular}{@{}>{\raggedright\arraybackslash}p{63.342mm}>{\raggedright\arraybackslash}p{57.846mm}@{}}
\toprule
\textbf{Person and number} & \textbf{Command} \\
\midrule
2 sg. & ādrīf \\
2 pl. & ādrīfað \\
\bottomrule
\end{tabular}
\par\endgroup
% course-panel:T15-P01:end
% mdtable:T15:end
% teaching: NEW-U12-GRAM-u12-fleon-imperative
% teaching: NEW-U12-PAR-fleon-IMP
\reserve{7\baselineskip}
% mdparagraph:P091:begin
\textbf{flēon; contracted strong verb: singular and plural commands}
% mdparagraph:P091:end

% table: NEW-U12-PAR-fleon-IMP
% mdtable:T16:begin
% course-table-style: 2026-10
% course-panel:T16-P01:begin
\par\begingroup
\fontsize{10.5}{13.2}\selectfont
\setlength{\LTleft}{0pt}\setlength{\LTright}{\fill}
\noindent
\begin{tabular}{@{}>{\raggedright\arraybackslash}p{63.342mm}>{\raggedright\arraybackslash}p{57.846mm}@{}}
\toprule
\textbf{Person and number} & \textbf{Command} \\
\midrule
2 sg. & flēoh \\
2 pl. & flēoð \\
\bottomrule
\end{tabular}
\par\endgroup
% course-panel:T16-P01:end
% mdtable:T16:end
% teaching: NEW-U12-GRAM-u12-gemetan-imperative
% teaching: NEW-U12-PAR-gemetan-IMP
\reserve{7\baselineskip}
% mdparagraph:P092:begin
\textbf{gemētan; weak verb: singular and plural commands}
% mdparagraph:P092:end

% table: NEW-U12-PAR-gemetan-IMP
% mdtable:T17:begin
% course-table-style: 2026-10
% course-panel:T17-P01:begin
\par\begingroup
\fontsize{10.5}{13.2}\selectfont
\setlength{\LTleft}{0pt}\setlength{\LTright}{\fill}
\noindent
\begin{tabular}{@{}>{\raggedright\arraybackslash}p{63.342mm}>{\raggedright\arraybackslash}p{57.846mm}@{}}
\toprule
\textbf{Person and number} & \textbf{Command} \\
\midrule
2 sg. & gemēt \\
2 pl. & gemētað \\
\bottomrule
\end{tabular}
\par\endgroup
% course-panel:T17-P01:end
% mdtable:T17:end
% teaching: NEW-U12-GRAM-u12-geopenian-imperative
% teaching: NEW-U12-PAR-geopenian-IMP
\reserve{7\baselineskip}
% mdparagraph:P093:begin
\textbf{geopenian; weak verb: singular and plural commands}
% mdparagraph:P093:end

% table: NEW-U12-PAR-geopenian-IMP
% mdtable:T18:begin
% course-table-style: 2026-10
% course-panel:T18-P01:begin
\par\begingroup
\fontsize{10.5}{13.2}\selectfont
\setlength{\LTleft}{0pt}\setlength{\LTright}{\fill}
\noindent
\begin{tabular}{@{}>{\raggedright\arraybackslash}p{63.342mm}>{\raggedright\arraybackslash}p{57.846mm}@{}}
\toprule
\textbf{Person and number} & \textbf{Command} \\
\midrule
2 sg. & geopena \\
2 pl. & geopeniað \\
\bottomrule
\end{tabular}
\par\endgroup
% course-panel:T18-P01:end
% mdtable:T18:end
% teaching: NEW-U12-GRAM-u12-weorpan-imperative
% teaching: NEW-U12-PAR-weorpan-IMP
\reserve{7\baselineskip}
% mdparagraph:P094:begin
\textbf{weorpan; strong verb: singular and plural commands}
% mdparagraph:P094:end

% table: NEW-U12-PAR-weorpan-IMP
% mdtable:T19:begin
% course-table-style: 2026-10
% course-panel:T19-P01:begin
\par\begingroup
\fontsize{10.5}{13.2}\selectfont
\setlength{\LTleft}{0pt}\setlength{\LTright}{\fill}
\noindent
\begin{tabular}{@{}>{\raggedright\arraybackslash}p{63.342mm}>{\raggedright\arraybackslash}p{57.846mm}@{}}
\toprule
\textbf{Person and number} & \textbf{Command} \\
\midrule
2 sg. & weorp \\
2 pl. & weorpað \\
\bottomrule
\end{tabular}
\par\endgroup
% course-panel:T19-P01:end
% mdtable:T19:end
% teaching: NEW-U12-GRAM-u12-leornian-imperative
% teaching: NEW-U12-PAR-leornian-IMP
\reserve{7\baselineskip}
% mdparagraph:P095:begin
\textbf{leornian; weak verb: singular and plural commands}
% mdparagraph:P095:end

% table: NEW-U12-PAR-leornian-IMP
% mdtable:T20:begin
% course-table-style: 2026-10
% course-panel:T20-P01:begin
\par\begingroup
\fontsize{10.5}{13.2}\selectfont
\setlength{\LTleft}{0pt}\setlength{\LTright}{\fill}
\noindent
\begin{tabular}{@{}>{\raggedright\arraybackslash}p{63.342mm}>{\raggedright\arraybackslash}p{57.846mm}@{}}
\toprule
\textbf{Person and number} & \textbf{Command} \\
\midrule
2 sg. & leorna \\
2 pl. & leorniað \\
\bottomrule
\end{tabular}
\par\endgroup
% course-panel:T20-P01:end
% mdtable:T20:end
% teaching: NEW-U12-GRAM-u12-risan-imperative
% teaching: NEW-U12-PAR-risan-IMP
\reserve{7\baselineskip}
% mdparagraph:P096:begin
\textbf{rīsan; strong verb: singular and plural commands}
% mdparagraph:P096:end

% table: NEW-U12-PAR-risan-IMP
% mdtable:T21:begin
% course-table-style: 2026-10
% course-panel:T21-P01:begin
\par\begingroup
\fontsize{10.5}{13.2}\selectfont
\setlength{\LTleft}{0pt}\setlength{\LTright}{\fill}
\noindent
\begin{tabular}{@{}>{\raggedright\arraybackslash}p{63.342mm}>{\raggedright\arraybackslash}p{57.846mm}@{}}
\toprule
\textbf{Person and number} & \textbf{Command} \\
\midrule
2 sg. & rīs \\
2 pl. & rīsað \\
\bottomrule
\end{tabular}
\par\endgroup
% course-panel:T21-P01:end
% mdtable:T21:end
% mdparagraph:P097:begin
The already learned \textbf{sleah} is the singular command of slēan. For rise use the newly displayed \textbf{rīs}; do not add a new prefix to the selected verb. Earlier \textbf{cum}, \textbf{gā}, \textbf{nim} and \textbf{sęle} remain available in their taught uses.
% mdparagraph:P097:end

\reserve{7\baselineskip}
% mdparagraph:P098:begin
\subsection*{Participle bases}
% mdparagraph:P098:end

% mdparagraph:P099:begin
A present participle characteristically has -ende; the past bases must be learned with their verbs. The first three pairs below are productive as base forms. Their uses remain limited to the constructions taught here or earlier. A base pair is not a complete adjectival declension.
% mdparagraph:P099:end

% teaching: NEW-U12-GRAM-u12-gewyrcean-participle-bases
% teaching: NEW-U12-PAR-gewyrcean-PART
\reserve{7\baselineskip}
% mdparagraph:P100:begin
\textbf{gewyrcean; weak verb: present and past participle bases}
% mdparagraph:P100:end

% table: NEW-U12-PAR-gewyrcean-PART
% mdtable:T22:begin
\par\begingroup
\fontsize{10.5}{13.2}\selectfont
\setlength{\LTleft}{0pt}\setlength{\LTright}{\fill}
\noindent
\begin{tabular}{@{}>{\raggedright\arraybackslash}p{69.84mm}>{\raggedright\arraybackslash}p{51.35mm}@{}}
\toprule
\textbf{Form} & \textbf{Base} \\
\midrule
Present participle & gewyrcende \\
Past participle & geworht \\
\bottomrule
\end{tabular}
\par\endgroup

% mdtable:T22:end
% teaching: NEW-U12-GRAM-u12-gierwan-participle-bases
% teaching: NEW-U12-PAR-gierwan-PART
\reserve{7\baselineskip}
% mdparagraph:P101:begin
\textbf{gierwan; weak verb: present and past participle bases}
% mdparagraph:P101:end

% table: NEW-U12-PAR-gierwan-PART
% mdtable:T23:begin
\par\begingroup
\fontsize{10.5}{13.2}\selectfont
\setlength{\LTleft}{0pt}\setlength{\LTright}{\fill}
\noindent
\begin{tabular}{@{}>{\raggedright\arraybackslash}p{69.84mm}>{\raggedright\arraybackslash}p{51.35mm}@{}}
\toprule
\textbf{Form} & \textbf{Base} \\
\midrule
Present participle & gierwende \\
Past participle & gierwed \\
\bottomrule
\end{tabular}
\par\endgroup

% mdtable:T23:end
% teaching: NEW-U12-GRAM-u12-settan-participle-bases
% teaching: NEW-U12-PAR-settan-PART
\reserve{7\baselineskip}
% mdparagraph:P102:begin
\textbf{sęttan; weak verb: present and past participle bases}
% mdparagraph:P102:end

% table: NEW-U12-PAR-settan-PART
% mdtable:T24:begin
\par\begingroup
\fontsize{10.5}{13.2}\selectfont
\setlength{\LTleft}{0pt}\setlength{\LTright}{\fill}
\noindent
\begin{tabular}{@{}>{\raggedright\arraybackslash}p{69.84mm}>{\raggedright\arraybackslash}p{51.35mm}@{}}
\toprule
\textbf{Form} & \textbf{Base} \\
\midrule
Present participle & sęttende \\
Past participle & sętt \\
\bottomrule
\end{tabular}
\par\endgroup

% mdtable:T24:end
% teaching: NEW-U12-GRAM-u12-gewitan-participle-bases
% mdparagraph:P103:begin
\textbf{Reference pairs.} Recognize both bases of gewītan and sittan. Only the stated constructions with gewiten and sittende are analyzed here; neither pair is assigned for independent composition.
% mdparagraph:P103:end

% teaching: NEW-U12-PAR-gewitan-PART
\reserve{7\baselineskip}
% mdparagraph:P104:begin
\textbf{gewītan; strong verb: present and past participle bases}
% mdparagraph:P104:end

% table: NEW-U12-PAR-gewitan-PART
% mdtable:T25:begin
\par\begingroup
\fontsize{10.5}{13.2}\selectfont
\setlength{\LTleft}{0pt}\setlength{\LTright}{\fill}
\noindent
\begin{tabular}{@{}>{\raggedright\arraybackslash}p{69.84mm}>{\raggedright\arraybackslash}p{51.35mm}@{}}
\toprule
\textbf{Form} & \textbf{Base} \\
\midrule
Present participle & gewītende \\
Past participle & gewiten \\
\bottomrule
\end{tabular}
\par\endgroup

% mdtable:T25:end
% teaching: NEW-U12-GRAM-u12-sittan-participle-bases
% teaching: NEW-U12-PAR-sittan-PART
\reserve{7\baselineskip}
% mdparagraph:P105:begin
\textbf{sittan; strong verb: present and past participle bases}
% mdparagraph:P105:end

% table: NEW-U12-PAR-sittan-PART
% mdtable:T26:begin
\par\begingroup
\fontsize{10.5}{13.2}\selectfont
\setlength{\LTleft}{0pt}\setlength{\LTright}{\fill}
\noindent
\begin{tabular}{@{}>{\raggedright\arraybackslash}p{69.84mm}>{\raggedright\arraybackslash}p{51.35mm}@{}}
\toprule
\textbf{Form} & \textbf{Base} \\
\midrule
Present participle & sittende \\
Past participle & seten \\
\bottomrule
\end{tabular}
\par\endgroup

% mdtable:T26:end
% mdparagraph:P106:begin
The exact selected bases \textbf{gierwed} and \textbf{sętt} stand beside the earlier citation spellings \textbf{(ge)gier(w)ed} and \textbf{(ge)sętt}. Do not treat the parentheses as letters to write in a sentence. The earlier supplied \textbf{geworht} is now available in the productive constructions of this unit.
% mdparagraph:P106:end

\reserve{7\baselineskip}
% mdparagraph:P107:begin
\subsection*{Plain and inflected infinitives}
% mdparagraph:P107:end

% mdparagraph:P108:begin
Use the plain infinitive after the assigned modal: ādrīfan after ability expressed by magan, and genēosian after intention expressed by willan. The inflected forms are shown for recognition only; their display does not add new composition with tō.
% mdparagraph:P108:end

% teaching: NEW-U12-GRAM-u12-adrifan-modal-infinitive
% teaching: NEW-U12-PAR-adrifan-INF
\reserve{7\baselineskip}
% mdparagraph:P109:begin
\textbf{ādrīfan; strong verb: plain and inflected infinitives}
% mdparagraph:P109:end

% table: NEW-U12-PAR-adrifan-INF
% mdtable:T27:begin
\par\begingroup
\fontsize{10.5}{13.2}\selectfont
\setlength{\LTleft}{0pt}\setlength{\LTright}{\fill}
\noindent
\begin{tabular}{@{}>{\raggedright\arraybackslash}p{71.72mm}>{\raggedright\arraybackslash}p{49.46mm}@{}}
\toprule
\textbf{Form} & \textbf{Infinitive} \\
\midrule
Plain infinitive & ādrīfan \\
Inflected infinitive & ādrīfanne \\
\bottomrule
\end{tabular}
\par\endgroup

% mdtable:T27:end
% teaching: NEW-U12-GRAM-u12-geneosian-modal-infinitive
% teaching: NEW-U12-PAR-geneosian-INF
\reserve{7\baselineskip}
% mdparagraph:P110:begin
\textbf{genēosian; weak verb: plain and inflected infinitives}
% mdparagraph:P110:end

% table: NEW-U12-PAR-geneosian-INF
% mdtable:T28:begin
\par\begingroup
\fontsize{10.5}{13.2}\selectfont
\setlength{\LTleft}{0pt}\setlength{\LTright}{\fill}
\noindent
\begin{tabular}{@{}>{\raggedright\arraybackslash}p{68.91mm}>{\raggedright\arraybackslash}p{52.28mm}@{}}
\toprule
\textbf{Form} & \textbf{Infinitive} \\
\midrule
Plain infinitive & genēosian \\
Inflected infinitive & genēosianne \\
\bottomrule
\end{tabular}
\par\endgroup

% mdtable:T28:end
\reserve{7\baselineskip}
% mdparagraph:P111:begin
\section*{6. Nouns: full patterns and lexical stems}
% mdparagraph:P111:end

% teaching: NEW-U12-NOMINAL-SYSTEMS
% teaching: NEW-U12-LEXICAL-DIAGNOSTICS
% teaching: NEW-U12-WEAK-NOMINALS
% mdparagraph:P112:begin
The complete patterns below give nominative, genitive, dative and accusative in both numbers. Where another noun uses the same pattern, its diagnostic row gives the headword, stem and distinguishing forms. The substitution instruction applies to \textbf{every corresponding case and number}, not merely to the examples in that row. Use the separate headword where expressly stated. Nouns keep their lexical gender.
% mdparagraph:P112:end

% mdparagraph:P113:begin
\textbf{Vowel change and syncope.} Fōt uses fēt in its selected dative singular and nominative/\allowbreak{}accusative plural. Fēond selects fēonde in the dative singular and fīend in the nominative/\allowbreak{}accusative plural; its genitive plural is fēonda. In ēðel, sāwol and wǣpen the medial vowel is absent from the ending-bearing stem: ēðl-, sāwl-, wǣpn-. Hærfest and pęning retain their medial vowel; the ordinary syncope rule requires a short second syllable.
% mdparagraph:P113:end

% mdparagraph:P114:begin
\textbf{Weak nouns.} Andwlita and sǣdtīma use the weak masculine pattern. Fǣmne and heorte use the weak feminine pattern, with fǣmn- and heort- before endings. The abstract suffix in gehīersumnes has -ness- before an ending. Gehīersumnes and ręst are productive in their assigned singular abstract meanings; their plural forms are \textbf{reference only}, without a new count-plural sense.
% mdparagraph:P114:end

% teaching: NEW-U12-VOC-BRYCG
% teaching: NEW-U12-NOM-BRYCG
% diagnostic: NEW-U12-DIAG-BRYCG
% Page review repair: keep the lexical lead with its first display or derivation.
\reserve{8\baselineskip}
% mdparagraph:P115:begin
\textbf{brycg} (feminine): bridge. Diagnostic stem: \textbf{brycg}.
% mdparagraph:P115:end

% mdtable:T29:begin
\par\begingroup
\fontsize{10.5}{13.2}\selectfont
\setlength{\LTleft}{0pt}\setlength{\LTright}{\fill}
\noindent
\begin{tabular}{@{}>{\raggedright\arraybackslash}p{39.24mm}>{\raggedright\arraybackslash}p{36.31mm}>{\raggedright\arraybackslash}p{42.82mm}@{}}
\toprule
\textbf{Genitive singular} & \textbf{Dative singular} & \textbf{Nominative plural} \\
\midrule
brycge & brycge & brycga \\
\bottomrule
\end{tabular}
\par\endgroup

% mdtable:T29:end
\reserve{9\baselineskip}
% mdparagraph:P116:begin
\textbf{brycg; feminine noun}
% mdparagraph:P116:end

% table: NEW-U12-NOM-BRYCG
% mdtable:T30:begin
\par\begingroup
\fontsize{10.5}{13.2}\selectfont
\setlength{\LTleft}{0pt}\setlength{\LTright}{\fill}
\noindent
\begin{tabular}{@{}>{\raggedright\arraybackslash}p{39.46mm}>{\raggedright\arraybackslash}p{39.46mm}>{\raggedright\arraybackslash}p{39.46mm}@{}}
\toprule
\textbf{Case} & \textbf{Singular} & \textbf{Plural} \\
\midrule
Nominative & brycg & brycga \\
Genitive & brycge & brycga \\
Dative & brycge & brycgum \\
Accusative & brycge & brycga \\
\bottomrule
\end{tabular}
\par\endgroup

% mdtable:T30:end
% mdparagraph:P117:begin
Selected nominative and accusative plural brycga have the accepted alternative \textbf{brycge}. The table uses -a consistently.
% mdparagraph:P117:end

% teaching: NEW-U12-VOC-GEHIERSUMNES
% teaching: NEW-U12-NOM-GEHIERSUMNES
% diagnostic: NEW-U12-DIAG-GEHIERSUMNES
% Page review repair: keep the lexical lead with its first display or derivation.
\reserve{8\baselineskip}
% mdparagraph:P118:begin
\textbf{gehīersumnes} (feminine): obedience. Diagnostic stem: \textbf{gehīersumness}.
% mdparagraph:P118:end

% mdtable:T31:begin
\par\begingroup
\fontsize{10.5}{13.2}\selectfont
\setlength{\LTleft}{0pt}\setlength{\LTright}{\fill}
\noindent
\begin{tabular}{@{}>{\raggedright\arraybackslash}p{40.25mm}>{\raggedright\arraybackslash}p{37.89mm}>{\raggedright\arraybackslash}p{40.25mm}@{}}
\toprule
\textbf{Genitive singular} & \textbf{Dative singular} & \textbf{Nominative plural} \\
\midrule
gehīersumnesse & gehīersumnesse & gehīersumnessa \\
\bottomrule
\end{tabular}
\par\endgroup

% mdtable:T31:end
\reserve{9\baselineskip}
% mdparagraph:P119:begin
\textbf{gehīersumnes; feminine noun}
% mdparagraph:P119:end

% table: NEW-U12-NOM-GEHIERSUMNES
% mdtable:T32:begin
\par\begingroup
\fontsize{10.5}{13.2}\selectfont
\setlength{\LTleft}{0pt}\setlength{\LTright}{\fill}
\noindent
\begin{tabular}{@{}>{\raggedright\arraybackslash}p{30.35mm}>{\raggedright\arraybackslash}p{42.49mm}>{\raggedright\arraybackslash}p{45.53mm}@{}}
\toprule
\textbf{Case} & \textbf{Singular} & \textbf{Plural} \\
\midrule
Nominative & gehīersumnes & gehīersumnessa \\
Genitive & gehīersumnesse & gehīersumnessa \\
Dative & gehīersumnesse & gehīersumnessum \\
Accusative & gehīersumnesse & gehīersumnessa \\
\bottomrule
\end{tabular}
\par\endgroup

% mdtable:T32:end
% mdparagraph:P120:begin
The entire plural column above is \textbf{reference only}. In the singular, obedience is an abstract thing promised.
% mdparagraph:P120:end

% teaching: NEW-U12-VOC-ANDWLITA
% teaching: NEW-U12-NOM-ANDWLITA
% diagnostic: NEW-U12-DIAG-ANDWLITA
% Page review repair: keep the lexical lead with its first display or derivation.
\reserve{8\baselineskip}
% mdparagraph:P121:begin
\textbf{andwlita} (masculine): face; countenance. Diagnostic stem: \textbf{andwlit}.
% mdparagraph:P121:end

% mdtable:T33:begin
\par\begingroup
\fontsize{10.5}{13.2}\selectfont
\setlength{\LTleft}{0pt}\setlength{\LTright}{\fill}
\noindent
\begin{tabular}{@{}>{\raggedright\arraybackslash}p{39.83mm}>{\raggedright\arraybackslash}p{36.98mm}>{\raggedright\arraybackslash}p{41.57mm}@{}}
\toprule
\textbf{Genitive singular} & \textbf{Dative singular} & \textbf{Nominative plural} \\
\midrule
andwlitan & andwlitan & andwlitan \\
\bottomrule
\end{tabular}
\par\endgroup

% mdtable:T33:end
\reserve{9\baselineskip}
% mdparagraph:P122:begin
\textbf{andwlita; masculine noun}
% mdparagraph:P122:end

% table: NEW-U12-NOM-ANDWLITA
% mdtable:T34:begin
\par\begingroup
\fontsize{10.5}{13.2}\selectfont
\setlength{\LTleft}{0pt}\setlength{\LTright}{\fill}
\noindent
\begin{tabular}{@{}>{\raggedright\arraybackslash}p{39.46mm}>{\raggedright\arraybackslash}p{39.46mm}>{\raggedright\arraybackslash}p{39.46mm}@{}}
\toprule
\textbf{Case} & \textbf{Singular} & \textbf{Plural} \\
\midrule
Nominative & andwlita & andwlitan \\
Genitive & andwlitan & andwlitena \\
Dative & andwlitan & andwlitum \\
Accusative & andwlitan & andwlitan \\
\bottomrule
\end{tabular}
\par\endgroup

% mdtable:T34:end
% teaching: NEW-U12-VOC-GEWRIT
% teaching: NEW-U12-NOM-GEWRIT
% diagnostic: NEW-U12-DIAG-GEWRIT
% Page review repair: keep the lexical lead with its first display or derivation.
\reserve{8\baselineskip}
% mdparagraph:P123:begin
\textbf{gewrit} (neuter): letter; written document. Diagnostic stem: \textbf{gewrit}.
% mdparagraph:P123:end

% mdtable:T35:begin
\par\begingroup
\fontsize{10.5}{13.2}\selectfont
\setlength{\LTleft}{0pt}\setlength{\LTright}{\fill}
\noindent
\begin{tabular}{@{}>{\raggedright\arraybackslash}p{39.24mm}>{\raggedright\arraybackslash}p{36.31mm}>{\raggedright\arraybackslash}p{42.82mm}@{}}
\toprule
\textbf{Genitive singular} & \textbf{Dative singular} & \textbf{Nominative plural} \\
\midrule
gewrites & gewrite & gewritu \\
\bottomrule
\end{tabular}
\par\endgroup

% mdtable:T35:end
\reserve{9\baselineskip}
% mdparagraph:P124:begin
\textbf{gewrit; neuter noun}
% mdparagraph:P124:end

% table: NEW-U12-NOM-GEWRIT
% mdtable:T36:begin
\par\begingroup
\fontsize{10.5}{13.2}\selectfont
\setlength{\LTleft}{0pt}\setlength{\LTright}{\fill}
\noindent
\begin{tabular}{@{}>{\raggedright\arraybackslash}p{39.46mm}>{\raggedright\arraybackslash}p{39.46mm}>{\raggedright\arraybackslash}p{39.46mm}@{}}
\toprule
\textbf{Case} & \textbf{Singular} & \textbf{Plural} \\
\midrule
Nominative & gewrit & gewritu \\
Genitive & gewrites & gewrita \\
Dative & gewrite & gewritum \\
Accusative & gewrit & gewritu \\
\bottomrule
\end{tabular}
\par\endgroup

% mdtable:T36:end
% teaching: NEW-U12-VOC-HERE
% teaching: NEW-U12-NOM-HERE
% diagnostic: NEW-U12-DIAG-HERE
% Page review repair: keep the lexical lead with its first display or derivation.
\reserve{8\baselineskip}
% mdparagraph:P125:begin
\textbf{hęre} (masculine): army. Diagnostic stem: \textbf{hęr}.
% mdparagraph:P125:end

% mdtable:T37:begin
\par\begingroup
\fontsize{10.5}{13.2}\selectfont
\setlength{\LTleft}{0pt}\setlength{\LTright}{\fill}
\noindent
\begin{tabular}{@{}>{\raggedright\arraybackslash}p{39.24mm}>{\raggedright\arraybackslash}p{36.31mm}>{\raggedright\arraybackslash}p{42.82mm}@{}}
\toprule
\textbf{Genitive singular} & \textbf{Dative singular} & \textbf{Nominative plural} \\
\midrule
hęres & hęre & hęras \\
\bottomrule
\end{tabular}
\par\endgroup

% mdtable:T37:end
\reserve{9\baselineskip}
% mdparagraph:P126:begin
\textbf{hęre; masculine noun}
% mdparagraph:P126:end

% table: NEW-U12-NOM-HERE
% mdtable:T38:begin
\par\begingroup
\fontsize{10.5}{13.2}\selectfont
\setlength{\LTleft}{0pt}\setlength{\LTright}{\fill}
\noindent
\begin{tabular}{@{}>{\raggedright\arraybackslash}p{39.46mm}>{\raggedright\arraybackslash}p{39.46mm}>{\raggedright\arraybackslash}p{39.46mm}@{}}
\toprule
\textbf{Case} & \textbf{Singular} & \textbf{Plural} \\
\midrule
Nominative & hęre & hęras \\
Genitive & hęres & hęra \\
Dative & hęre & hęrum \\
Accusative & hęre & hęras \\
\bottomrule
\end{tabular}
\par\endgroup

% mdtable:T38:end
% mdparagraph:P127:begin
The selected ordinary stem is hęr-. Cook also allows \textbf{hęrg-} or \textbf{hęrig-} before singular endings and \textbf{hęrg-, hęrig-} or \textbf{hęrige-} before plural endings. Thus hęres can have hęrges/\allowbreak{}hęriges; hęras can have hęrgas/\allowbreak{}hęrigas/\allowbreak{}hęrigeas. Apply each stem before the corresponding ending, including genitive and dative endings. These attested alternatives are accepted, not additional recall targets.
% mdparagraph:P127:end

% teaching: NEW-U12-VOC-SCIPHERE
% teaching: NEW-U12-NOM-SCIPHERE
% diagnostic: NEW-U12-DIAG-SCIPHERE
% Page review repair: keep the lexical lead with its first display or derivation.
\reserve{8\baselineskip}
% mdparagraph:P128:begin
\textbf{sciphęre} (masculine): fleet; naval force. Diagnostic stem: \textbf{sciphęr}.
% mdparagraph:P128:end

% mdtable:T39:begin
\par\begingroup
\fontsize{10.5}{13.2}\selectfont
\setlength{\LTleft}{0pt}\setlength{\LTright}{\fill}
\noindent
\begin{tabular}{@{}>{\raggedright\arraybackslash}p{40.42mm}>{\raggedright\arraybackslash}p{35.77mm}>{\raggedright\arraybackslash}p{42.19mm}@{}}
\toprule
\textbf{Genitive singular} & \textbf{Dative singular} & \textbf{Nominative plural} \\
\midrule
sciphęres & sciphęre & sciphęras \\
\bottomrule
\end{tabular}
\par\endgroup

% mdtable:T39:end
% mdparagraph:P129:begin
Use the complete \textbf{hęre; masculine noun} table. In every corresponding cell replace initial \textbf{hęr-} with \textbf{sciphęr-}, retaining the whole printed ending and the same case and number. An empty ending remains empty.
% mdparagraph:P129:end

% mdparagraph:P130:begin
The same alternatives occur in the compound: sciphęrg-/\allowbreak{}sciphęrig- before singular endings and sciphęrg-/\allowbreak{}sciphęrig-/\allowbreak{}sciphęrige- before plural endings. The selected complete model remains sciphęr- + the hęre endings.
% mdparagraph:P130:end

% teaching: NEW-U12-VOC-FEOND
% teaching: NEW-U12-NOM-FEOND
% diagnostic: NEW-U12-DIAG-FEOND
% Page review repair: keep the lexical lead with its first display or derivation.
\reserve{8\baselineskip}
% mdparagraph:P131:begin
\textbf{fēond} (masculine): enemy; foe. Diagnostic stem: \textbf{fēond /\allowbreak{} fīend}.
% mdparagraph:P131:end

% mdtable:T40:begin
\par\begingroup
\fontsize{10.5}{13.2}\selectfont
\setlength{\LTleft}{0pt}\setlength{\LTright}{\fill}
\noindent
\begin{tabular}{@{}>{\raggedright\arraybackslash}p{39.24mm}>{\raggedright\arraybackslash}p{36.31mm}>{\raggedright\arraybackslash}p{42.82mm}@{}}
\toprule
\textbf{Genitive singular} & \textbf{Dative singular} & \textbf{Nominative plural} \\
\midrule
fēondes & fēonde & fīend \\
\bottomrule
\end{tabular}
\par\endgroup

% mdtable:T40:end
\reserve{9\baselineskip}
% mdparagraph:P132:begin
\textbf{fēond; masculine noun}
% mdparagraph:P132:end

% table: NEW-U12-NOM-FEOND
% mdtable:T41:begin
\par\begingroup
\fontsize{10.5}{13.2}\selectfont
\setlength{\LTleft}{0pt}\setlength{\LTright}{\fill}
\noindent
\begin{tabular}{@{}>{\raggedright\arraybackslash}p{39.46mm}>{\raggedright\arraybackslash}p{39.46mm}>{\raggedright\arraybackslash}p{39.46mm}@{}}
\toprule
\textbf{Case} & \textbf{Singular} & \textbf{Plural} \\
\midrule
Nominative & fēond & fīend \\
Genitive & fēondes & fēonda \\
Dative & fēonde & fēondum \\
Accusative & fēond & fīend \\
\bottomrule
\end{tabular}
\par\endgroup

% mdtable:T41:end
% mdparagraph:P133:begin
Accepted alternatives are \textbf{fīend} in the dative singular and \textbf{fēond} or \textbf{fēondas} in the nominative/\allowbreak{}accusative plural. Learn selected genitive plural \textbf{fēonda}; it does not take the -endra of the warrior family.
% mdparagraph:P133:end

% teaching: NEW-U12-VOC-HAERFEST
% teaching: NEW-U12-NOM-HAERFEST
% diagnostic: NEW-U12-DIAG-HAERFEST
% Page review repair: keep the lexical lead with its first display or derivation.
\reserve{8\baselineskip}
% mdparagraph:P134:begin
\textbf{hærfest} (masculine): harvest. Diagnostic stem: \textbf{hærfest}.
% mdparagraph:P134:end

% mdtable:T42:begin
\par\begingroup
\fontsize{10.5}{13.2}\selectfont
\setlength{\LTleft}{0pt}\setlength{\LTright}{\fill}
\noindent
\begin{tabular}{@{}>{\raggedright\arraybackslash}p{40.42mm}>{\raggedright\arraybackslash}p{35.77mm}>{\raggedright\arraybackslash}p{42.19mm}@{}}
\toprule
\textbf{Genitive singular} & \textbf{Dative singular} & \textbf{Nominative plural} \\
\midrule
hærfestes & hærfeste & hærfestas \\
\bottomrule
\end{tabular}
\par\endgroup

% mdtable:T42:end
\reserve{9\baselineskip}
% mdparagraph:P135:begin
\textbf{hærfest; masculine noun}
% mdparagraph:P135:end

% table: NEW-U12-NOM-HAERFEST
% mdtable:T43:begin
\par\begingroup
\fontsize{10.5}{13.2}\selectfont
\setlength{\LTleft}{0pt}\setlength{\LTright}{\fill}
\noindent
\begin{tabular}{@{}>{\raggedright\arraybackslash}p{39.46mm}>{\raggedright\arraybackslash}p{39.46mm}>{\raggedright\arraybackslash}p{39.46mm}@{}}
\toprule
\textbf{Case} & \textbf{Singular} & \textbf{Plural} \\
\midrule
Nominative & hærfest & hærfestas \\
Genitive & hærfestes & hærfesta \\
Dative & hærfeste & hærfestum \\
Accusative & hærfest & hærfestas \\
\bottomrule
\end{tabular}
\par\endgroup

% mdtable:T43:end
% teaching: NEW-U12-VOC-SAEDTIMA
% teaching: NEW-U12-NOM-SAEDTIMA
% diagnostic: NEW-U12-DIAG-SAEDTIMA
% Page review repair: keep the lexical lead with its first display or derivation.
\reserve{8\baselineskip}
% mdparagraph:P136:begin
\textbf{sǣdtīma} (masculine): seedtime; sowing season. Diagnostic stem: \textbf{sǣdtīm}.
% mdparagraph:P136:end

% mdtable:T44:begin
\par\begingroup
\fontsize{10.5}{13.2}\selectfont
\setlength{\LTleft}{0pt}\setlength{\LTright}{\fill}
\noindent
\begin{tabular}{@{}>{\raggedright\arraybackslash}p{39.24mm}>{\raggedright\arraybackslash}p{36.31mm}>{\raggedright\arraybackslash}p{42.82mm}@{}}
\toprule
\textbf{Genitive singular} & \textbf{Dative singular} & \textbf{Nominative plural} \\
\midrule
sǣdtīman & sǣdtīman & sǣdtīman \\
\bottomrule
\end{tabular}
\par\endgroup

% mdtable:T44:end
% mdparagraph:P137:begin
Use the complete \textbf{andwlita; masculine noun} table. In every corresponding cell replace initial \textbf{andwlit-} with \textbf{sǣdtīm-}, retaining the whole printed ending and the same case and number. An empty ending remains empty.
% mdparagraph:P137:end

% teaching: NEW-U12-VOC-WAEPEN
% teaching: NEW-U12-NOM-WAEPEN
% diagnostic: NEW-U12-DIAG-WAEPEN
% Page review repair: keep the lexical lead with its first display or derivation.
\reserve{8\baselineskip}
% mdparagraph:P138:begin
\textbf{wǣpen} (neuter): weapon. Diagnostic stem: \textbf{wǣpn}.
% mdparagraph:P138:end

% mdtable:T45:begin
\par\begingroup
\fontsize{10.5}{13.2}\selectfont
\setlength{\LTleft}{0pt}\setlength{\LTright}{\fill}
\noindent
\begin{tabular}{@{}>{\raggedright\arraybackslash}p{39.24mm}>{\raggedright\arraybackslash}p{36.31mm}>{\raggedright\arraybackslash}p{42.82mm}@{}}
\toprule
\textbf{Genitive singular} & \textbf{Dative singular} & \textbf{Nominative plural} \\
\midrule
wǣpnes & wǣpne & wǣpnu \\
\bottomrule
\end{tabular}
\par\endgroup

% mdtable:T45:end
\reserve{9\baselineskip}
% mdparagraph:P139:begin
\textbf{wǣpen; neuter noun}
% mdparagraph:P139:end

% table: NEW-U12-NOM-WAEPEN
% mdtable:T46:begin
\par\begingroup
\fontsize{10.5}{13.2}\selectfont
\setlength{\LTleft}{0pt}\setlength{\LTright}{\fill}
\noindent
\begin{tabular}{@{}>{\raggedright\arraybackslash}p{39.46mm}>{\raggedright\arraybackslash}p{39.46mm}>{\raggedright\arraybackslash}p{39.46mm}@{}}
\toprule
\textbf{Case} & \textbf{Singular} & \textbf{Plural} \\
\midrule
Nominative & wǣpen & wǣpnu \\
Genitive & wǣpnes & wǣpna \\
Dative & wǣpne & wǣpnum \\
Accusative & wǣpen & wǣpnu \\
\bottomrule
\end{tabular}
\par\endgroup

% mdtable:T46:end
% mdparagraph:P140:begin
The selected nominative/\allowbreak{}accusative plural is \textbf{wǣpnu}. Cook also records \textbf{wǣpen} and \textbf{wǣpn} in these positions. Retain the shortened stem wǣpn- before the other endings.
% mdparagraph:P140:end

% teaching: NEW-U12-VOC-WULF
% teaching: NEW-U12-NOM-WULF
% diagnostic: NEW-U12-DIAG-WULF
% Page review repair: keep the lexical lead with its first display or derivation.
\reserve{8\baselineskip}
% mdparagraph:P141:begin
\textbf{wulf} (masculine): wolf. Diagnostic stem: \textbf{wulf}.
% mdparagraph:P141:end

% mdtable:T47:begin
\par\begingroup
\fontsize{10.5}{13.2}\selectfont
\setlength{\LTleft}{0pt}\setlength{\LTright}{\fill}
\noindent
\begin{tabular}{@{}>{\raggedright\arraybackslash}p{39.24mm}>{\raggedright\arraybackslash}p{36.31mm}>{\raggedright\arraybackslash}p{42.82mm}@{}}
\toprule
\textbf{Genitive singular} & \textbf{Dative singular} & \textbf{Nominative plural} \\
\midrule
wulfes & wulfe & wulfas \\
\bottomrule
\end{tabular}
\par\endgroup

% mdtable:T47:end
% mdparagraph:P142:begin
Use the complete \textbf{hærfest; masculine noun} table. In every corresponding cell replace initial \textbf{hærfest-} with \textbf{wulf-}, retaining the whole printed ending and the same case and number. An empty ending remains empty.
% mdparagraph:P142:end

% teaching: NEW-U12-VOC-GEMAERE
% teaching: NEW-U12-NOM-GEMAERE
% diagnostic: NEW-U12-DIAG-GEMAERE
% Page review repair: keep the lexical lead with its first display or derivation.
\reserve{8\baselineskip}
% mdparagraph:P143:begin
\textbf{gemǣre} (neuter): boundary; border. Diagnostic stem: \textbf{gemǣr}.
% mdparagraph:P143:end

% mdtable:T48:begin
\par\begingroup
\fontsize{10.5}{13.2}\selectfont
\setlength{\LTleft}{0pt}\setlength{\LTright}{\fill}
\noindent
\begin{tabular}{@{}>{\raggedright\arraybackslash}p{39.24mm}>{\raggedright\arraybackslash}p{36.31mm}>{\raggedright\arraybackslash}p{42.82mm}@{}}
\toprule
\textbf{Genitive singular} & \textbf{Dative singular} & \textbf{Nominative plural} \\
\midrule
gemǣres & gemǣre & gemǣru \\
\bottomrule
\end{tabular}
\par\endgroup

% mdtable:T48:end
\reserve{9\baselineskip}
% mdparagraph:P144:begin
\textbf{gemǣre; neuter noun}
% mdparagraph:P144:end

% table: NEW-U12-NOM-GEMAERE
% mdtable:T49:begin
\par\begingroup
\fontsize{10.5}{13.2}\selectfont
\setlength{\LTleft}{0pt}\setlength{\LTright}{\fill}
\noindent
\begin{tabular}{@{}>{\raggedright\arraybackslash}p{39.46mm}>{\raggedright\arraybackslash}p{39.46mm}>{\raggedright\arraybackslash}p{39.46mm}@{}}
\toprule
\textbf{Case} & \textbf{Singular} & \textbf{Plural} \\
\midrule
Nominative & gemǣre & gemǣru \\
Genitive & gemǣres & gemǣra \\
Dative & gemǣre & gemǣrum \\
Accusative & gemǣre & gemǣru \\
\bottomrule
\end{tabular}
\par\endgroup

% mdtable:T49:end
% teaching: NEW-U12-VOC-PENING
% teaching: NEW-U12-NOM-PENING
% diagnostic: NEW-U12-DIAG-PENING
% Page review repair: keep the lexical lead with its first display or derivation.
\reserve{8\baselineskip}
% mdparagraph:P145:begin
\textbf{pęning} (masculine): penny; coin. Diagnostic stem: \textbf{pęning}.
% mdparagraph:P145:end

% mdtable:T50:begin
\par\begingroup
\fontsize{10.5}{13.2}\selectfont
\setlength{\LTleft}{0pt}\setlength{\LTright}{\fill}
\noindent
\begin{tabular}{@{}>{\raggedright\arraybackslash}p{39.24mm}>{\raggedright\arraybackslash}p{36.31mm}>{\raggedright\arraybackslash}p{42.82mm}@{}}
\toprule
\textbf{Genitive singular} & \textbf{Dative singular} & \textbf{Nominative plural} \\
\midrule
pęninges & pęninge & pęningas \\
\bottomrule
\end{tabular}
\par\endgroup

% mdtable:T50:end
% mdparagraph:P146:begin
Use the complete \textbf{hærfest; masculine noun} table. In every corresponding cell replace initial \textbf{hærfest-} with \textbf{pęning-}, retaining the whole printed ending and the same case and number. An empty ending remains empty.
% mdparagraph:P146:end

% teaching: NEW-U12-VOC-MUTH
% teaching: NEW-U12-NOM-MUTH
% diagnostic: NEW-U12-DIAG-MUTH
% Page review repair: keep the lexical lead with its first display or derivation.
\reserve{8\baselineskip}
% mdparagraph:P147:begin
\textbf{mūð} (masculine): mouth. Diagnostic stem: \textbf{mūð}.
% mdparagraph:P147:end

% mdtable:T51:begin
\par\begingroup
\fontsize{10.5}{13.2}\selectfont
\setlength{\LTleft}{0pt}\setlength{\LTright}{\fill}
\noindent
\begin{tabular}{@{}>{\raggedright\arraybackslash}p{39.24mm}>{\raggedright\arraybackslash}p{36.31mm}>{\raggedright\arraybackslash}p{42.82mm}@{}}
\toprule
\textbf{Genitive singular} & \textbf{Dative singular} & \textbf{Nominative plural} \\
\midrule
mūðes & mūðe & mūðas \\
\bottomrule
\end{tabular}
\par\endgroup

% mdtable:T51:end
% mdparagraph:P148:begin
Use the complete \textbf{hærfest; masculine noun} table. In every corresponding cell replace initial \textbf{hærfest-} with \textbf{mūð-}, retaining the whole printed ending and the same case and number. An empty ending remains empty.
% mdparagraph:P148:end

% teaching: NEW-U12-VOC-FAEMNE
% teaching: NEW-U12-NOM-FAEMNE
% diagnostic: NEW-U12-DIAG-FAEMNE
% Page review repair: keep the lexical lead with its first display or derivation.
\reserve{8\baselineskip}
% mdparagraph:P149:begin
\textbf{fǣmne} (feminine): maiden; young woman. Diagnostic stem: \textbf{fǣmn}.
% mdparagraph:P149:end

% mdtable:T52:begin
\par\begingroup
\fontsize{10.5}{13.2}\selectfont
\setlength{\LTleft}{0pt}\setlength{\LTright}{\fill}
\noindent
\begin{tabular}{@{}>{\raggedright\arraybackslash}p{39.24mm}>{\raggedright\arraybackslash}p{36.31mm}>{\raggedright\arraybackslash}p{42.82mm}@{}}
\toprule
\textbf{Genitive singular} & \textbf{Dative singular} & \textbf{Nominative plural} \\
\midrule
fǣmnan & fǣmnan & fǣmnan \\
\bottomrule
\end{tabular}
\par\endgroup

% mdtable:T52:end
\reserve{9\baselineskip}
% mdparagraph:P150:begin
\textbf{fǣmne; feminine noun}
% mdparagraph:P150:end

% table: NEW-U12-NOM-FAEMNE
% mdtable:T53:begin
\par\begingroup
\fontsize{10.5}{13.2}\selectfont
\setlength{\LTleft}{0pt}\setlength{\LTright}{\fill}
\noindent
\begin{tabular}{@{}>{\raggedright\arraybackslash}p{39.46mm}>{\raggedright\arraybackslash}p{39.46mm}>{\raggedright\arraybackslash}p{39.46mm}@{}}
\toprule
\textbf{Case} & \textbf{Singular} & \textbf{Plural} \\
\midrule
Nominative & fǣmne & fǣmnan \\
Genitive & fǣmnan & fǣmnena \\
Dative & fǣmnan & fǣmnum \\
Accusative & fǣmnan & fǣmnan \\
\bottomrule
\end{tabular}
\par\endgroup

% mdtable:T53:end
% teaching: NEW-U12-VOC-STANWEALL
% teaching: NEW-U12-NOM-STANWEALL
% diagnostic: NEW-U12-DIAG-STANWEALL
% Page review repair: keep the lexical lead with its first display or derivation.
\reserve{8\baselineskip}
% mdparagraph:P151:begin
\textbf{stānweall} (masculine): stone wall. Diagnostic stem: \textbf{stānweall}.
% mdparagraph:P151:end

% mdtable:T54:begin
\par\begingroup
\fontsize{10.5}{13.2}\selectfont
\setlength{\LTleft}{0pt}\setlength{\LTright}{\fill}
\noindent
\begin{tabular}{@{}>{\raggedright\arraybackslash}p{40.9mm}>{\raggedright\arraybackslash}p{36.57mm}>{\raggedright\arraybackslash}p{40.9mm}@{}}
\toprule
\textbf{Genitive singular} & \textbf{Dative singular} & \textbf{Nominative plural} \\
\midrule
stānwealles & stānwealle & stānweallas \\
\bottomrule
\end{tabular}
\par\endgroup

% mdtable:T54:end
% mdparagraph:P152:begin
Use the complete \textbf{hærfest; masculine noun} table. In every corresponding cell replace initial \textbf{hærfest-} with \textbf{stānweall-}, retaining the whole printed ending and the same case and number. An empty ending remains empty.
% mdparagraph:P152:end

% teaching: NEW-U12-VOC-LAND
% teaching: NEW-U12-NOM-LAND
% diagnostic: NEW-U12-DIAG-LAND
% Page review repair: keep the lexical lead with its first display or derivation.
\reserve{8\baselineskip}
% mdparagraph:P153:begin
\textbf{land} (neuter): land; country. Diagnostic stem: \textbf{land}.
% mdparagraph:P153:end

% mdtable:T55:begin
\par\begingroup
\fontsize{10.5}{13.2}\selectfont
\setlength{\LTleft}{0pt}\setlength{\LTright}{\fill}
\noindent
\begin{tabular}{@{}>{\raggedright\arraybackslash}p{39.24mm}>{\raggedright\arraybackslash}p{36.31mm}>{\raggedright\arraybackslash}p{42.82mm}@{}}
\toprule
\textbf{Genitive singular} & \textbf{Dative singular} & \textbf{Nominative plural} \\
\midrule
landes & lande & land \\
\bottomrule
\end{tabular}
\par\endgroup

% mdtable:T55:end
\reserve{9\baselineskip}
% mdparagraph:P154:begin
\textbf{land; neuter noun}
% mdparagraph:P154:end

% table: NEW-U12-NOM-LAND
% mdtable:T56:begin
\par\begingroup
\fontsize{10.5}{13.2}\selectfont
\setlength{\LTleft}{0pt}\setlength{\LTright}{\fill}
\noindent
\begin{tabular}{@{}>{\raggedright\arraybackslash}p{39.46mm}>{\raggedright\arraybackslash}p{39.46mm}>{\raggedright\arraybackslash}p{39.46mm}@{}}
\toprule
\textbf{Case} & \textbf{Singular} & \textbf{Plural} \\
\midrule
Nominative & land & land \\
Genitive & landes & landa \\
Dative & lande & landum \\
Accusative & land & land \\
\bottomrule
\end{tabular}
\par\endgroup

% mdtable:T56:end
% teaching: NEW-U12-VOC-FOT
% teaching: NEW-U12-NOM-FOT
% diagnostic: NEW-U12-DIAG-FOT
% Page review repair: keep the lexical lead with its first display or derivation.
\reserve{8\baselineskip}
% mdparagraph:P155:begin
\textbf{fōt} (masculine): foot; also a foot as a measure of length. Diagnostic stem: \textbf{fōt /\allowbreak{} fēt}.
% mdparagraph:P155:end

% mdtable:T57:begin
\par\begingroup
\fontsize{10.5}{13.2}\selectfont
\setlength{\LTleft}{0pt}\setlength{\LTright}{\fill}
\noindent
\begin{tabular}{@{}>{\raggedright\arraybackslash}p{39.24mm}>{\raggedright\arraybackslash}p{36.31mm}>{\raggedright\arraybackslash}p{42.82mm}@{}}
\toprule
\textbf{Genitive singular} & \textbf{Dative singular} & \textbf{Nominative plural} \\
\midrule
fōtes & fēt & fēt \\
\bottomrule
\end{tabular}
\par\endgroup

% mdtable:T57:end
\reserve{9\baselineskip}
% mdparagraph:P156:begin
\textbf{fōt; masculine noun}
% mdparagraph:P156:end

% table: NEW-U12-NOM-FOT
% mdtable:T58:begin
\par\begingroup
\fontsize{10.5}{13.2}\selectfont
\setlength{\LTleft}{0pt}\setlength{\LTright}{\fill}
\noindent
\begin{tabular}{@{}>{\raggedright\arraybackslash}p{39.46mm}>{\raggedright\arraybackslash}p{39.46mm}>{\raggedright\arraybackslash}p{39.46mm}@{}}
\toprule
\textbf{Case} & \textbf{Singular} & \textbf{Plural} \\
\midrule
Nominative & fōt & fēt \\
Genitive & fōtes & fōta \\
Dative & fēt & fōtum \\
Accusative & fōt & fēt \\
\bottomrule
\end{tabular}
\par\endgroup

% mdtable:T58:end
% mdparagraph:P157:begin
Accepted dative singular \textbf{fōte} and occasional nominative/\allowbreak{}accusative plural \textbf{fōtas} occur beside the selected fēt. Use \textbf{fōta} for the genitive plural measure expression taught in §8.
% mdparagraph:P157:end

% teaching: NEW-U12-VOC-ETHEL
% teaching: NEW-U12-NOM-ETHEL
% diagnostic: NEW-U12-DIAG-ETHEL
% Page review repair: keep the lexical lead with its first display or derivation.
\reserve{8\baselineskip}
% mdparagraph:P158:begin
\textbf{ēðel} (masculine): native land; homeland. Diagnostic stem: \textbf{ēðl}.
% mdparagraph:P158:end

% mdtable:T59:begin
\par\begingroup
\fontsize{10.5}{13.2}\selectfont
\setlength{\LTleft}{0pt}\setlength{\LTright}{\fill}
\noindent
\begin{tabular}{@{}>{\raggedright\arraybackslash}p{39.24mm}>{\raggedright\arraybackslash}p{36.31mm}>{\raggedright\arraybackslash}p{42.82mm}@{}}
\toprule
\textbf{Genitive singular} & \textbf{Dative singular} & \textbf{Nominative plural} \\
\midrule
ēðles & ēðle & ēðlas \\
\bottomrule
\end{tabular}
\par\endgroup

% mdtable:T59:end
\reserve{9\baselineskip}
% mdparagraph:P159:begin
\textbf{ēðel; masculine noun}
% mdparagraph:P159:end

% table: NEW-U12-NOM-ETHEL
% mdtable:T60:begin
\par\begingroup
\fontsize{10.5}{13.2}\selectfont
\setlength{\LTleft}{0pt}\setlength{\LTright}{\fill}
\noindent
\begin{tabular}{@{}>{\raggedright\arraybackslash}p{39.46mm}>{\raggedright\arraybackslash}p{39.46mm}>{\raggedright\arraybackslash}p{39.46mm}@{}}
\toprule
\textbf{Case} & \textbf{Singular} & \textbf{Plural} \\
\midrule
Nominative & ēðel & ēðlas \\
Genitive & ēðles & ēðla \\
Dative & ēðle & ēðlum \\
Accusative & ēðel & ēðlas \\
\bottomrule
\end{tabular}
\par\endgroup

% mdtable:T60:end
% teaching: NEW-U12-VOC-HEORTE
% teaching: NEW-U12-NOM-HEORTE
% diagnostic: NEW-U12-DIAG-HEORTE
% Page review repair: keep the lexical lead with its first display or derivation.
\reserve{8\baselineskip}
% mdparagraph:P160:begin
\textbf{heorte} (feminine): heart. Diagnostic stem: \textbf{heort}.
% mdparagraph:P160:end

% mdtable:T61:begin
\par\begingroup
\fontsize{10.5}{13.2}\selectfont
\setlength{\LTleft}{0pt}\setlength{\LTright}{\fill}
\noindent
\begin{tabular}{@{}>{\raggedright\arraybackslash}p{39.24mm}>{\raggedright\arraybackslash}p{36.31mm}>{\raggedright\arraybackslash}p{42.82mm}@{}}
\toprule
\textbf{Genitive singular} & \textbf{Dative singular} & \textbf{Nominative plural} \\
\midrule
heortan & heortan & heortan \\
\bottomrule
\end{tabular}
\par\endgroup

% mdtable:T61:end
% mdparagraph:P161:begin
Use the complete \textbf{fǣmne; feminine noun} table. In every corresponding cell replace initial \textbf{fǣmn-} with \textbf{heort-}, retaining the whole printed ending and the same case and number. An empty ending remains empty.
% mdparagraph:P161:end

% teaching: NEW-U12-VOC-REST
% teaching: NEW-U12-NOM-REST
% diagnostic: NEW-U12-DIAG-REST
% Page review repair: keep the lexical lead with its first display or derivation.
\reserve{8\baselineskip}
% mdparagraph:P162:begin
\textbf{ręst} (feminine): rest; repose. Diagnostic stem: \textbf{ręst}.
% mdparagraph:P162:end

% mdtable:T62:begin
\par\begingroup
\fontsize{10.5}{13.2}\selectfont
\setlength{\LTleft}{0pt}\setlength{\LTright}{\fill}
\noindent
\begin{tabular}{@{}>{\raggedright\arraybackslash}p{39.24mm}>{\raggedright\arraybackslash}p{36.31mm}>{\raggedright\arraybackslash}p{42.82mm}@{}}
\toprule
\textbf{Genitive singular} & \textbf{Dative singular} & \textbf{Nominative plural} \\
\midrule
ręste & ręste & ręsta \\
\bottomrule
\end{tabular}
\par\endgroup

% mdtable:T62:end
% mdparagraph:P163:begin
Use the complete \textbf{brycg; feminine noun} table. In every corresponding cell replace initial \textbf{brycg-} with \textbf{ręst-}, retaining the whole printed ending and the same case and number. An empty ending remains empty.
% mdparagraph:P163:end

% mdparagraph:P164:begin
All derived plural forms of ręst are \textbf{reference only}. The assigned productive sense is singular “rest; repose”, supported by Cook VI.21.9. It is not restricted to the vocabulary’s separate “couch; bed” gloss.
% mdparagraph:P164:end

% teaching: NEW-U12-VOC-SAWOL
% teaching: NEW-U12-NOM-SAWOL
% diagnostic: NEW-U12-DIAG-SAWOL
% Page review repair: keep the lexical lead with its first display or derivation.
\reserve{8\baselineskip}
% mdparagraph:P165:begin
\textbf{sāwol} (feminine): soul. Diagnostic stem: \textbf{sāwl}.
% mdparagraph:P165:end

% mdtable:T63:begin
\par\begingroup
\fontsize{10.5}{13.2}\selectfont
\setlength{\LTleft}{0pt}\setlength{\LTright}{\fill}
\noindent
\begin{tabular}{@{}>{\raggedright\arraybackslash}p{39.24mm}>{\raggedright\arraybackslash}p{36.31mm}>{\raggedright\arraybackslash}p{42.82mm}@{}}
\toprule
\textbf{Genitive singular} & \textbf{Dative singular} & \textbf{Nominative plural} \\
\midrule
sāwle & sāwle & sāwla \\
\bottomrule
\end{tabular}
\par\endgroup

% mdtable:T63:end
\reserve{9\baselineskip}
% mdparagraph:P166:begin
\textbf{sāwol; feminine noun}
% mdparagraph:P166:end

% table: NEW-U12-NOM-SAWOL
% mdtable:T64:begin
\par\begingroup
\fontsize{10.5}{13.2}\selectfont
\setlength{\LTleft}{0pt}\setlength{\LTright}{\fill}
\noindent
\begin{tabular}{@{}>{\raggedright\arraybackslash}p{39.46mm}>{\raggedright\arraybackslash}p{39.46mm}>{\raggedright\arraybackslash}p{39.46mm}@{}}
\toprule
\textbf{Case} & \textbf{Singular} & \textbf{Plural} \\
\midrule
Nominative & sāwol & sāwla \\
Genitive & sāwle & sāwla \\
Dative & sāwle & sāwlum \\
Accusative & sāwle & sāwla \\
\bottomrule
\end{tabular}
\par\endgroup

% mdtable:T64:end
% mdparagraph:P167:begin
Cook also gives the citation form \textbf{sāwl}. The selected headword is sāwol; ending-bearing forms use sāwl-.
% mdparagraph:P167:end

\reserve{7\baselineskip}
% mdparagraph:P168:begin
\section*{7. Adjectives: complete strong and weak patterns}
% mdparagraph:P168:end

% teaching: NEW-U12-ADJECTIVE-DIAGNOSTICS
% mdparagraph:P169:begin
An ordinary positive predicate adjective is strong and nominative, agreeing with its subject. An attributive adjective is weak after the article or demonstrative, strong in the ordinary bare phrase. Agree with the noun, not with the possessor. In direct address the noun has its nominative shape and an accompanying adjective is weak.
% mdparagraph:P169:end

% mdparagraph:P170:begin
Brād, beald, earg, ēaðmōd and mildheort use the complete long strong model and the weak model. Ēadig and the adjective īsern use the selected disyllabic short-model endings. Each model below contains all three genders in both numbers. Distinct instrumental singular cells are \textbf{reference only}.
% mdparagraph:P170:end

% teaching: NEW-U12-VOC-BRAD
% teaching: NEW-U12-ADJ-BRAD-STRONG
% teaching: NEW-U12-ADJ-BRAD-WEAK
% diagnostic: NEW-U12-DIAG-BRAD
% Page review repair: keep the lexical lead with its first display or derivation.
\reserve{24\baselineskip}
% mdparagraph:P171:begin
\textbf{brād} (adjective): broad. Diagnostic stem: \textbf{brād}.
% mdparagraph:P171:end

\reserve{18\baselineskip}
% mdparagraph:P172:begin
\textbf{brād; strong adjective}
% mdparagraph:P172:end

% table: NEW-U12-ADJ-BRAD-STRONG
% mdtable:T65:begin
% course-table-style: 2026-10
% course-panel:T65-P01:begin
\par\begingroup
\fontsize{10.5}{13.2}\selectfont
\setlength{\LTleft}{0pt}\setlength{\LTright}{\fill}
\noindent
\begin{tabular}{@{}>{\raggedright\arraybackslash}p{22.0mm}>{\raggedright\arraybackslash}p{31.188mm}>{\raggedright\arraybackslash}p{31.188mm}>{\raggedright\arraybackslash}p{31.188mm}@{}}
\toprule
\textbf{Case} & \textbf{Masculine singular} & \textbf{Neuter singular} & \textbf{Feminine singular} \\
\midrule
Nominative & brād & brād & brād \\
Genitive & brādes & brādes & brādre \\
Dative & brādum & brādum & brādre \\
Accusative & brādne & brād & brāde \\
Instrumental & brāde & brāde & as dative \\
\bottomrule
\end{tabular}
\par\endgroup
% course-panel:T65-P01:end
\par\addvspace{5pt}
% course-panel:T65-P02:begin
\par\begingroup
\fontsize{10.5}{13.2}\selectfont
\setlength{\LTleft}{0pt}\setlength{\LTright}{\fill}
\noindent
\begin{tabular}{@{}>{\raggedright\arraybackslash}p{22.0mm}>{\raggedright\arraybackslash}p{31.188mm}>{\raggedright\arraybackslash}p{31.188mm}>{\raggedright\arraybackslash}p{31.188mm}@{}}
\toprule
\textbf{Case} & \textbf{Masculine plural} & \textbf{Neuter plural} & \textbf{Feminine plural} \\
\midrule
Nominative & brāde & brād & brāda \\
Genitive & brādra & brādra & brādra \\
Dative & brādum & brādum & brādum \\
Accusative & brāde & brād & brāda \\
Instrumental & as dative & as dative & as dative \\
\bottomrule
\end{tabular}
\par\endgroup
% course-panel:T65-P02:end
% mdtable:T65:end

\reserve{18\baselineskip}
% mdparagraph:P173:begin
\textbf{brād; weak adjective}
% mdparagraph:P173:end

% table: NEW-U12-ADJ-BRAD-WEAK
% mdtable:T66:begin
% course-table-style: 2026-10
% course-panel:T66-P01:begin
\par\begingroup
\fontsize{10.5}{13.2}\selectfont
\setlength{\LTleft}{0pt}\setlength{\LTright}{\fill}
\noindent
\begin{tabular}{@{}>{\raggedright\arraybackslash}p{22.0mm}>{\raggedright\arraybackslash}p{31.188mm}>{\raggedright\arraybackslash}p{31.188mm}>{\raggedright\arraybackslash}p{31.188mm}@{}}
\toprule
\textbf{Case} & \textbf{Masculine singular} & \textbf{Neuter singular} & \textbf{Feminine singular} \\
\midrule
Nominative & brāda & brāde & brāde \\
Genitive & brādan & brādan & brādan \\
Dative & brādan & brādan & brādan \\
Accusative & brādan & brāde & brādan \\
Instrumental & brādan & brādan & as dative \\
\bottomrule
\end{tabular}
\par\endgroup
% course-panel:T66-P01:end
\par\addvspace{5pt}
% course-panel:T66-P02:begin
\par\begingroup
\fontsize{10.5}{13.2}\selectfont
\setlength{\LTleft}{0pt}\setlength{\LTright}{\fill}
\noindent
\begin{tabular}{@{}>{\raggedright\arraybackslash}p{22.0mm}>{\raggedright\arraybackslash}p{31.188mm}>{\raggedright\arraybackslash}p{31.188mm}>{\raggedright\arraybackslash}p{31.188mm}@{}}
\toprule
\textbf{Case} & \textbf{Masculine plural} & \textbf{Neuter plural} & \textbf{Feminine plural} \\
\midrule
Nominative & brādan & brādan & brādan \\
Genitive & brādra & brādra & brādra \\
Dative & brādum & brādum & brādum \\
Accusative & brādan & brādan & brādan \\
Instrumental & as dative & as dative & as dative \\
\bottomrule
\end{tabular}
\par\endgroup
% course-panel:T66-P02:end
% mdtable:T66:end

% teaching: NEW-U12-VOC-BEALD
% teaching: NEW-U12-ADJ-BEALD-STRONG
% teaching: NEW-U12-ADJ-BEALD-WEAK
% diagnostic: NEW-U12-DIAG-BEALD
% Page review repair: keep the lexical lead with its first display or derivation.
\reserve{10\baselineskip}
% mdparagraph:P174:begin
\textbf{beald} (adjective): bold; brave. Diagnostic stem: \textbf{beald}.
% mdparagraph:P174:end

% mdparagraph:P175:begin
Use the complete \textbf{brād; strong adjective} table. In every corresponding cell replace initial \textbf{brād-} with \textbf{beald-}, retaining the whole printed ending and the same case, number, gender and declension. An empty ending remains empty.
% mdparagraph:P175:end

% mdparagraph:P176:begin
Use the complete \textbf{brād; weak adjective} table. In every corresponding cell replace initial \textbf{brād-} with \textbf{beald-}, retaining the whole printed ending and the same case, number, gender and declension. An empty ending remains empty.
% mdparagraph:P176:end

% teaching: NEW-U12-VOC-EARG
% teaching: NEW-U12-ADJ-EARG-STRONG
% teaching: NEW-U12-ADJ-EARG-WEAK
% diagnostic: NEW-U12-DIAG-EARG
% Page review repair: keep the lexical lead with its first display or derivation.
\reserve{10\baselineskip}
% mdparagraph:P177:begin
\textbf{earg} (adjective): cowardly. Diagnostic stem: \textbf{earg}.
% mdparagraph:P177:end

% mdparagraph:P178:begin
Use the complete \textbf{brād; strong adjective} table. In every corresponding cell replace initial \textbf{brād-} with \textbf{earg-}, retaining the whole printed ending and the same case, number, gender and declension. An empty ending remains empty.
% mdparagraph:P178:end

% mdparagraph:P179:begin
Use the complete \textbf{brād; weak adjective} table. In every corresponding cell replace initial \textbf{brād-} with \textbf{earg-}, retaining the whole printed ending and the same case, number, gender and declension. An empty ending remains empty.
% mdparagraph:P179:end

% teaching: NEW-U12-VOC-EADMOD
% teaching: NEW-U12-ADJ-EADMOD-STRONG
% teaching: NEW-U12-ADJ-EADMOD-WEAK
% diagnostic: NEW-U12-DIAG-EADMOD
% Page review repair: keep the lexical lead with its first display or derivation.
\reserve{10\baselineskip}
% mdparagraph:P180:begin
\textbf{ēaðmōd} (adjective): lowly; humble. Diagnostic stem: \textbf{ēaðmōd}.
% mdparagraph:P180:end

% mdparagraph:P181:begin
Use the complete \textbf{brād; strong adjective} table. In every corresponding cell replace initial \textbf{brād-} with \textbf{ēaðmōd-}, retaining the whole printed ending and the same case, number, gender and declension. An empty ending remains empty.
% mdparagraph:P181:end

% mdparagraph:P182:begin
Use the complete \textbf{brād; weak adjective} table. In every corresponding cell replace initial \textbf{brād-} with \textbf{ēaðmōd-}, retaining the whole printed ending and the same case, number, gender and declension. An empty ending remains empty.
% mdparagraph:P182:end

% teaching: NEW-U12-VOC-EADIG
% teaching: NEW-U12-ADJ-EADIG-STRONG
% teaching: NEW-U12-ADJ-EADIG-WEAK
% diagnostic: NEW-U12-DIAG-EADIG
% Page review repair: keep the lexical lead with its first display or derivation.
\reserve{24\baselineskip}
% mdparagraph:P183:begin
\textbf{ēadig} (adjective): happy; blessed. Diagnostic stem: \textbf{ēadig}.
% mdparagraph:P183:end

\reserve{18\baselineskip}
% mdparagraph:P184:begin
\textbf{ēadig; strong adjective}
% mdparagraph:P184:end

% table: NEW-U12-ADJ-EADIG-STRONG
% mdtable:T67:begin
% course-table-style: 2026-10
% course-panel:T67-P01:begin
\par\begingroup
\fontsize{10.5}{13.2}\selectfont
\setlength{\LTleft}{0pt}\setlength{\LTright}{\fill}
\noindent
\begin{tabular}{@{}>{\raggedright\arraybackslash}p{22.0mm}>{\raggedright\arraybackslash}p{31.188mm}>{\raggedright\arraybackslash}p{31.188mm}>{\raggedright\arraybackslash}p{31.188mm}@{}}
\toprule
\textbf{Case} & \textbf{Masculine singular} & \textbf{Neuter singular} & \textbf{Feminine singular} \\
\midrule
Nominative & ēadig & ēadig & ēadigu \\
Genitive & ēadiges & ēadiges & ēadigre \\
Dative & ēadigum & ēadigum & ēadigre \\
Accusative & ēadigne & ēadig & ēadige \\
Instrumental & ēadige & ēadige & as dative \\
\bottomrule
\end{tabular}
\par\endgroup
% course-panel:T67-P01:end
\par\addvspace{5pt}
% course-panel:T67-P02:begin
\par\begingroup
\fontsize{10.5}{13.2}\selectfont
\setlength{\LTleft}{0pt}\setlength{\LTright}{\fill}
\noindent
\begin{tabular}{@{}>{\raggedright\arraybackslash}p{22.0mm}>{\raggedright\arraybackslash}p{31.188mm}>{\raggedright\arraybackslash}p{31.188mm}>{\raggedright\arraybackslash}p{31.188mm}@{}}
\toprule
\textbf{Case} & \textbf{Masculine plural} & \textbf{Neuter plural} & \textbf{Feminine plural} \\
\midrule
Nominative & ēadige & ēadigu & ēadiga \\
Genitive & ēadigra & ēadigra & ēadigra \\
Dative & ēadigum & ēadigum & ēadigum \\
Accusative & ēadige & ēadigu & ēadiga \\
Instrumental & as dative & as dative & as dative \\
\bottomrule
\end{tabular}
\par\endgroup
% course-panel:T67-P02:end
% mdtable:T67:end

\reserve{18\baselineskip}
% mdparagraph:P185:begin
\textbf{ēadig; weak adjective}
% mdparagraph:P185:end

% table: NEW-U12-ADJ-EADIG-WEAK
% mdtable:T68:begin
% course-table-style: 2026-10
% course-panel:T68-P01:begin
\par\begingroup
\fontsize{10.5}{13.2}\selectfont
\setlength{\LTleft}{0pt}\setlength{\LTright}{\fill}
\noindent
\begin{tabular}{@{}>{\raggedright\arraybackslash}p{22.0mm}>{\raggedright\arraybackslash}p{31.188mm}>{\raggedright\arraybackslash}p{31.188mm}>{\raggedright\arraybackslash}p{31.188mm}@{}}
\toprule
\textbf{Case} & \textbf{Masculine singular} & \textbf{Neuter singular} & \textbf{Feminine singular} \\
\midrule
Nominative & ēadiga & ēadige & ēadige \\
Genitive & ēadigan & ēadigan & ēadigan \\
Dative & ēadigan & ēadigan & ēadigan \\
Accusative & ēadigan & ēadige & ēadigan \\
Instrumental & ēadigan & ēadigan & as dative \\
\bottomrule
\end{tabular}
\par\endgroup
% course-panel:T68-P01:end
\par\addvspace{5pt}
% course-panel:T68-P02:begin
\par\begingroup
\fontsize{10.5}{13.2}\selectfont
\setlength{\LTleft}{0pt}\setlength{\LTright}{\fill}
\noindent
\begin{tabular}{@{}>{\raggedright\arraybackslash}p{22.0mm}>{\raggedright\arraybackslash}p{31.188mm}>{\raggedright\arraybackslash}p{31.188mm}>{\raggedright\arraybackslash}p{31.188mm}@{}}
\toprule
\textbf{Case} & \textbf{Masculine plural} & \textbf{Neuter plural} & \textbf{Feminine plural} \\
\midrule
Nominative & ēadigan & ēadigan & ēadigan \\
Genitive & ēadigra & ēadigra & ēadigra \\
Dative & ēadigum & ēadigum & ēadigum \\
Accusative & ēadigan & ēadigan & ēadigan \\
Instrumental & as dative & as dative & as dative \\
\bottomrule
\end{tabular}
\par\endgroup
% course-panel:T68-P02:end
% mdtable:T68:end

% mdparagraph:P186:begin
The selected table retains i throughout. Cook also permits syncope before vowel endings, as in \textbf{ēadges}, and endingless feminine nominative singular and neuter nominative/\allowbreak{}accusative plural alternatives. These evidenced alternatives are accepted; they are not additional required tables.
% mdparagraph:P186:end

% teaching: NEW-U12-VOC-MILDHEORT
% teaching: NEW-U12-ADJ-MILDHEORT-STRONG
% teaching: NEW-U12-ADJ-MILDHEORT-WEAK
% diagnostic: NEW-U12-DIAG-MILDHEORT
% Page review repair: keep the lexical lead with its first display or derivation.
\reserve{10\baselineskip}
% mdparagraph:P187:begin
\textbf{mildheort} (adjective): merciful; compassionate. Diagnostic stem: \textbf{mildheort}.
% mdparagraph:P187:end

% mdparagraph:P188:begin
Use the complete \textbf{brād; strong adjective} table. In every corresponding cell replace initial \textbf{brād-} with \textbf{mildheort-}, retaining the whole printed ending and the same case, number, gender and declension. An empty ending remains empty.
% mdparagraph:P188:end

% mdparagraph:P189:begin
Use the complete \textbf{brād; weak adjective} table. In every corresponding cell replace initial \textbf{brād-} with \textbf{mildheort-}, retaining the whole printed ending and the same case, number, gender and declension. An empty ending remains empty.
% mdparagraph:P189:end

% teaching: NEW-U12-VOC-ISERN-ADJECTIVE
% teaching: NEW-U12-ADJ-ISERN-ADJECTIVE-STRONG
% teaching: NEW-U12-ADJ-ISERN-ADJECTIVE-WEAK
% diagnostic: NEW-U12-DIAG-ISERN-ADJECTIVE
% Page review repair: keep the lexical lead with its first display or derivation.
\reserve{10\baselineskip}
% mdparagraph:P190:begin
\textbf{īsern} (adjective): made of iron. Diagnostic stem: \textbf{īsern}.
% mdparagraph:P190:end

% mdparagraph:P191:begin
Use the complete \textbf{ēadig; strong adjective} table. In every corresponding cell replace initial \textbf{ēadig-} with \textbf{īsern-}, retaining the whole printed ending and the same case, number, gender and declension. An empty ending remains empty.
% mdparagraph:P191:end

% mdparagraph:P192:begin
Use the complete \textbf{ēadig; weak adjective} table. In every corresponding cell replace initial \textbf{ēadig-} with \textbf{īsern-}, retaining the whole printed ending and the same case, number, gender and declension. An empty ending remains empty.
% mdparagraph:P192:end

% mdparagraph:P193:begin
This \textbf{īsern} means “made of iron” and agrees as an adjective. It is distinct from the already learned noun meaning iron. Retain its medial e in the selected forms.
% mdparagraph:P193:end

\reserve{7\baselineskip}
% mdparagraph:P194:begin
\section*{8. Additional words and constructions}
% mdparagraph:P194:end

\reserve{7\baselineskip}
% mdparagraph:P195:begin
\subsection*{Adverbs}
% mdparagraph:P195:end

% teaching: NEW-U12-VOC-HAM
% teaching: NEW-U12-VOC-HIDER
% teaching: NEW-U12-VOC-GIET
% teaching: NEW-U12-VOC-TO-DAEG
% teaching: NEW-U12-VOC-NAEFRE
% mdtable:T69:begin
\par\begingroup
\fontsize{10.5}{13.2}\selectfont
\setlength{\LTleft}{0pt}\setlength{\LTright}{\fill}
\noindent
\begin{tabular}{@{}>{\raggedright\arraybackslash}p{28.45mm}>{\raggedright\arraybackslash}p{92.74mm}@{}}
\toprule
\textbf{Word} & \textbf{Meaning and use} \\
\midrule
hām & home; homewards, expressing direction without a preposition \\
hider & hither; toward the speaker \\
gīet & yet; still \\
tō-dæg & today \\
nǣfre & never \\
\bottomrule
\end{tabular}
\par\endgroup

% mdtable:T69:end
\reserve{7\baselineskip}
% mdparagraph:P196:begin
\subsection*{Objects and recipients}
% mdparagraph:P196:end

% teaching: NEW-U12-GRAM-u12-gierwan-accusative
% teaching: NEW-U12-GRAM-u12-writan-accusative
% teaching: NEW-U12-GRAM-u12-sendan-recipient
% teaching: NEW-U12-GRAM-u12-weorpan-accusative
% teaching: NEW-U12-GRAM-u12-gemetan-beneficiary
% mdparagraph:P197:begin
\textbf{Gierwan} takes the accusative thing prepared. The same thing remains accusative with habban + gierwed. \textbf{Wrītan} takes an accusative text or letter written. \textbf{Sęndan} takes the accusative person or thing sent and, if expressed, a dative recipient. \textbf{Weorpan} takes the accusative thing thrown, including a fishhook. Distinguish \textbf{angel}, fishhook, from \textbf{ęngel}, angel.
% mdparagraph:P197:end

% mdparagraph:P198:begin
With \textbf{gemētan}, the thing found is accusative. In the assigned expression about finding rest, the souls benefiting from it are dative. An English “unto” does not require a preposition there. This is the stated beneficiary use, not a rule assigning dative to every object of finding.
% mdparagraph:P198:end

% teaching: NEW-U12-GRAM-u12-geswican-intransitive
% teaching: NEW-U12-GRAM-u12-slean-kill
% teaching: NEW-U12-GRAM-u12-geseon-objectless
% teaching: NEW-U12-VOC-GESWICAN
% teaching: NEW-U12-VOC-SLEAN
% teaching: NEW-U12-VOC-GESEON
% mdparagraph:P199:begin
The earlier \textbf{geswīcan + genitive} expresses ceasing from an activity or rule. Also use geswīcan \textbf{without an expressed object} when seedtime and harvest cease or fail. \textbf{Slēan} can mean kill or slay, with an accusative person; this adds a lethal sense to strike/\allowbreak{}smite. \textbf{Gesēon} can stand without an expressed object in the assigned personification of the sea. Keep the earlier selected masculine gender of \textbf{sǣ} when choosing its pronoun.
% mdparagraph:P199:end

\reserve{7\baselineskip}
% mdparagraph:P200:begin
\subsection*{Prepositions and the meanings of their phrases}
% mdparagraph:P200:end

% teaching: NEW-U12-GRAM-u12-on-respect
% teaching: NEW-U12-GRAM-u12-welig-on-dative
% teaching: NEW-U12-GRAM-u12-beforan-motion
% teaching: NEW-U12-GRAM-u12-abutan-accusative
% teaching: NEW-U12-GRAM-u12-gemang-accusative
% teaching: NEW-U12-GRAM-u12-aet-learning-source
% teaching: NEW-U12-GRAM-u12-for-beneficiary
% teaching: NEW-U12-VOC-ON
% teaching: NEW-U12-VOC-WELIG
% teaching: NEW-U12-VOC-BEFORAN
% teaching: NEW-U12-VOC-ABUTAN
% teaching: NEW-U12-VOC-GEMANG
% teaching: NEW-U12-VOC-AET
% teaching: NEW-U12-VOC-FOR
% mdtable:T70:begin
\par\begingroup
\fontsize{10.5}{13.2}\selectfont
\setlength{\LTleft}{0pt}\setlength{\LTright}{\fill}
\noindent
\begin{tabular}{@{}>{\raggedright\arraybackslash}p{34.65mm}>{\raggedright\arraybackslash}p{86.54mm}@{}}
\toprule
\textbf{Construction} & \textbf{Use} \\
\midrule
welig on + dative & rich in a resource or product; this can describe a harvest as well as a person or place \\
ēaðmōd on + dative & lowly in a stated respect; on heortan, in heart \\
sęndan … beforan + dative & send someone before/\allowbreak{}ahead of a person; distinct from a dative recipient \\
ābūtan + accusative & around a person or thing, including building around it \\
sęndan … gemang + accusative & send among a group; use the accusative in the assigned sheep-and-wolves expression \\
leornian æt + dative & learn from a person; the thing learned may remain unexpressed \\
sęllan for + dative & give for/\allowbreak{}on behalf of a person in the assigned payment expression \\
\bottomrule
\end{tabular}
\par\endgroup

% mdtable:T70:end
% mdparagraph:P201:begin
The first two on phrases express respect, not movement. Beforan with the dative extends the earlier static “before” phrase to the person ahead of whom someone is sent. Dative location after gemang also occurs, but the productive motion phrase assigned here uses accusative. The for construction here does not teach every case or causal sense of that preposition.
% mdparagraph:P201:end

\reserve{7\baselineskip}
% mdparagraph:P202:begin
\subsection*{Duration, extent, and possession}
% mdparagraph:P202:end

% teaching: NEW-U12-GRAM-u12-duration-accusative
% teaching: NEW-U12-GRAM-u12-extent-genitive
% teaching: NEW-U12-GRAM-u12-predicate-possession
% mdparagraph:P203:begin
An accusative quantity of time can express \textbf{duration}: how long an action or state lasts. This differs from identifying the day on which it happens. Higher numerals may govern a genitive noun; the phrase for forty years is supplied at Cook 19.2, so do not force its noun into an accusative shape.
% mdparagraph:P203:end

% mdparagraph:P204:begin
A numeral with a \textbf{genitive measure noun} can state breadth or height. The noun fōt has genitive plural \textbf{fōta}. The number of feet measures the extent; the predicate adjective still agrees with the subject being described. The twenty-feet phrase is supplied at Cook 21.3.
% mdparagraph:P204:end

% mdparagraph:P205:begin
Possession can also be predicative. A \textbf{genitive possessor} after be means “is someone’s”; the genitive is not a nominative noun predicate. An \textbf{agreeing possessive adjective} after be agrees with what is possessed. Thus the person/\allowbreak{}number of the possessor selects ðīn, while the gender and number of the possessed dæg or niht determine its grammatical position. The same surface can occupy more than one such position.
% mdparagraph:P205:end

\reserve{7\baselineskip}
% mdparagraph:P206:begin
\subsection*{Comparison}
% mdparagraph:P206:end

% teaching: NEW-U12-GRAM-u12-more-numerous
% teaching: NEW-U12-GRAM-u12-equative-swa-swa
% teaching: NEW-U12-GRAM-u12-swa-predicative-comparison
% teaching: NEW-U12-VOC-MA
% teaching: NEW-U12-VOC-SWA
% mdparagraph:P207:begin
With a plural subject, indeclinable \textbf{mā} can mean “more in number”. A comparison standard after ðonne uses the earlier parallel-subject construction. This quantitative predicate is distinct from the earlier use of mā for a greater amount of doing; it does not require a new comparative adjective māran.
% mdparagraph:P207:end

% mdparagraph:P208:begin
For equality of degree, use \textbf{swā … adjective … swā}. The first swā marks the degree, “as”; the second introduces the standard, “as”. A positive predicate adjective retains strong agreement with the subject. The nominal standard can be nominative, understood as the subject of be. These are two separate words with an adjective between them, not the hyphenated source conjunction swā-swā.
% mdparagraph:P208:end

% mdparagraph:P209:begin
A single \textbf{swā + noun phrase} can instead compare the persons sent with sheep: “as; like”. The comparing noun phrase refers to the persons who are the object of sending. This is the particular nominal comparison assigned here, distinct from the earlier comparison between finite clauses.
% mdparagraph:P209:end

\reserve{7\baselineskip}
% mdparagraph:P210:begin
\subsection*{Never and never again}
% mdparagraph:P210:end

% teaching: NEW-U12-GRAM-u12-negative-concord-bounded
% mdparagraph:P211:begin
In the assigned \textbf{nǣfre … ne} construction, nǣfre and the negative finite clause reinforce one negative meaning. Place ne before the finite verb. \textbf{Nǣfre mā} means “never more; never again”. Do not cancel the two negatives as though they made an affirmative. Some verbs have contracted negative forms; nis, nāh and the required nolde remain visibly supplied below. This introduces the stated expressions only, not an entire family of negative pronouns or negative paradigms.
% mdparagraph:P211:end

\reserve{7\baselineskip}
% mdparagraph:P212:begin
\section*{9. Practice}
% mdparagraph:P212:end

% mdparagraph:P213:begin
Work through the form exercises before the analyses and compositions. Use the complete printed systems, preserving quantity and stem changes. In analysis, state both the form and its function: a shared spelling alone may not decide mood, case or person. Reference constructions may be analyzed with their explanation in view.
% mdparagraph:P213:end

\reserve{7\baselineskip}
% mdparagraph:P214:begin
\subsection*{F01}
% mdparagraph:P214:end

\par\begin{minipage}{\linewidth}\raggedright
% item: NEW-U12-F01
% mdparagraph:P215:begin
Give the complete nominative, genitive, dative and accusative singular and plural of fōt. Explain the vowel alternation.
% mdparagraph:P215:end

\end{minipage}\par\addvspace{4pt}

\reserve{7\baselineskip}
% mdparagraph:P216:begin
\subsection*{F02}
% mdparagraph:P216:end

\par\begin{minipage}{\linewidth}\raggedright
% item: NEW-U12-F02
% mdparagraph:P217:begin
Give the complete nominative, genitive, dative and accusative singular and plural of brycg and hęre, using the selected tables. Explain the differing stems before endings.
% mdparagraph:P217:end

\end{minipage}\par\addvspace{4pt}

\reserve{7\baselineskip}
% mdparagraph:P218:begin
\subsection*{F03}
% mdparagraph:P218:end

\par\begin{minipage}{\linewidth}\raggedright
% item: NEW-U12-F03
% mdparagraph:P219:begin
Give the complete present and preterite indicative of flēon and grētan: all three persons in both numbers for each tense.
% mdparagraph:P219:end

\end{minipage}\par\addvspace{4pt}

\reserve{7\baselineskip}
% mdparagraph:P220:begin
\subsection*{F04}
% mdparagraph:P220:end

\par\begin{minipage}{\linewidth}\raggedright
% item: NEW-U12-F04
% mdparagraph:P221:begin
For gemētan, ādrīfan, geopenian and cȳðan, give the present indicative second and third singular, the preterite second singular and the preterite third plural. Explain the stem and ending choices.
% mdparagraph:P221:end

\end{minipage}\par\addvspace{4pt}

\reserve{7\baselineskip}
% mdparagraph:P222:begin
\subsection*{F05}
% mdparagraph:P222:end

\par\begin{minipage}{\linewidth}\raggedright
% item: NEW-U12-F05
% mdparagraph:P223:begin
Give the singular and plural commands of flēon, ādrīfan, gemētan and geopenian.
% mdparagraph:P223:end

\end{minipage}\par\addvspace{4pt}

\reserve{7\baselineskip}
% mdparagraph:P224:begin
\subsection*{F06}
% mdparagraph:P224:end

\par\begin{minipage}{\linewidth}\raggedright
% item: NEW-U12-F06
% mdparagraph:P225:begin
Give all six person-and-number forms of the present indicative of habban with ðæt hūs geworht, then all six preterite indicative forms with the same words. Label the finite auxiliary and the unchanged past participle.
% mdparagraph:P225:end

\end{minipage}\par\addvspace{4pt}

\reserve{7\baselineskip}
% mdparagraph:P226:begin
\subsection*{F07}
% mdparagraph:P226:end

\par\begin{minipage}{\linewidth}\raggedright
% item: NEW-U12-F07
% mdparagraph:P227:begin
Give the base present and past participles of gierwan, sęttan and gewyrcean. Explain which base is used in the habban construction.
% mdparagraph:P227:end

\end{minipage}\par\addvspace{4pt}

\reserve{7\baselineskip}
% mdparagraph:P228:begin
\subsection*{A01}
% mdparagraph:P228:end

\par\begin{minipage}{\linewidth}\raggedright
% item: NEW-U12-A01
% mdparagraph:P229:begin
Translate and put the two narrated events in order. Identify the finite verb in each clause.
% mdparagraph:P229:end

% mdparagraph:P230:begin
Se mann cōm, and hē geseah ðæt hūs.
% mdparagraph:P230:end

\end{minipage}\par\addvspace{4pt}

\reserve{7\baselineskip}
% mdparagraph:P231:begin
\subsection*{A02}
% mdparagraph:P231:end

\par\begin{minipage}{\linewidth}\raggedright
% item: NEW-U12-A02
% mdparagraph:P232:begin
The house was finished before the visit. Translate, preserving that order of events. Identify the morphological tense of geworhte and explain why English may use had built.
% mdparagraph:P232:end

% mdparagraph:P233:begin
Se mann geseah ðæt hūs ðe his sunu geworhte.
% mdparagraph:P233:end

\end{minipage}\par\addvspace{4pt}

\reserve{7\baselineskip}
% mdparagraph:P234:begin
\subsection*{A03}
% mdparagraph:P234:end

\par\begin{minipage}{\linewidth}\raggedright
% item: NEW-U12-A03
% mdparagraph:P235:begin
The speaker is describing completed building work that they performed. Translate both sentences. Identify the builder, the object, the finite auxiliary and the past participle. Explain the difference between the two temporal viewpoints.
% mdparagraph:P235:end

% mdparagraph:P236:begin
Ic hæbbe ðæt hūs geworht.
% mdparagraph:P236:end

% mdparagraph:P237:begin
Ic hæfde ðæt hūs geworht.
% mdparagraph:P237:end

\end{minipage}\par\addvspace{4pt}

\reserve{7\baselineskip}
% mdparagraph:P238:begin
\subsection*{A04}
% mdparagraph:P238:end

\par\begin{minipage}{\linewidth}\raggedright
% item: NEW-U12-A04
% mdparagraph:P239:begin
In the second sentence, the speaker has just finished building the house. Translate both sentences. What work does habban do in each, and what additional information is introduced by geworht?
% mdparagraph:P239:end

% mdparagraph:P240:begin
Ic hæbbe ðæt hūs.
% mdparagraph:P240:end

% mdparagraph:P241:begin
Ic hæbbe ðæt hūs geworht.
% mdparagraph:P241:end

\end{minipage}\par\addvspace{4pt}

\reserve{7\baselineskip}
% mdparagraph:P242:begin
\subsection*{A05}
% mdparagraph:P242:end

\par\begin{minipage}{\linewidth}\raggedright
% item: NEW-U12-A05
% mdparagraph:P243:begin
Translate all three sentences. Identify their subjects and the horse’s role. Which sentence names the binder? Explain the usual event/\allowbreak{}state distinction without treating it as an absolute rule.
% mdparagraph:P243:end

% mdparagraph:P244:begin
Se mann band ðæt hors.
% mdparagraph:P244:end

% mdparagraph:P245:begin
Ðæt hors wearð bunden.
% mdparagraph:P245:end

% mdparagraph:P246:begin
Ðæt hors wæs bunden.
% mdparagraph:P246:end

\end{minipage}\par\addvspace{4pt}

\reserve{7\baselineskip}
% mdparagraph:P247:begin
\subsection*{A06}
% mdparagraph:P247:end

\par\begin{minipage}{\linewidth}\raggedright
% item: NEW-U12-A06
% mdparagraph:P248:begin
Translate. Identify the finite verb and the intransitive past participle. Explain why the man is not the object of an unnamed person’s action.
% mdparagraph:P248:end

% mdparagraph:P249:begin
Se mann wæs gewiten.
% mdparagraph:P249:end

\end{minipage}\par\addvspace{4pt}

\reserve{7\baselineskip}
% mdparagraph:P250:begin
\subsection*{A07}
% mdparagraph:P250:end

\par\begin{minipage}{\linewidth}\raggedright
% item: NEW-U12-A07
% mdparagraph:P251:begin
Translate, giving a simple-past rendering and a contextually possible English progressive. Identify the two verb forms and explain why the English construction is not imposed by form alone.
% mdparagraph:P251:end

% mdparagraph:P252:begin
Se mann wæs sittende.
% mdparagraph:P252:end

\end{minipage}\par\addvspace{4pt}

\reserve{7\baselineskip}
% mdparagraph:P253:begin
\subsection*{A08}
% mdparagraph:P253:end

\par\begin{minipage}{\linewidth}\raggedright
% item: NEW-U12-A08
% mdparagraph:P254:begin
Translate both reports. In each, the reported presence is simultaneous with the saying. Identify the main-clause tense and dependent mood; explain which moment the reported presence is related to.
% mdparagraph:P254:end

% mdparagraph:P255:begin
Hē sægð ðæt hē sīe hēr.
% mdparagraph:P255:end

% mdparagraph:P256:begin
Hē sægde ðæt hē wǣre hēr.
% mdparagraph:P256:end

\end{minipage}\par\addvspace{4pt}

\reserve{7\baselineskip}
% mdparagraph:P257:begin
\subsection*{A09}
% mdparagraph:P257:end

\par\begin{minipage}{\linewidth}\raggedright
% item: NEW-U12-A09
% mdparagraph:P258:begin
Translate. In this report the building precedes the saying. Identify the finite verbs, the participle and the object, and explain how the temporal relation is expressed.
% mdparagraph:P258:end

% mdparagraph:P259:begin
Hē sægde ðæt hē ðæt hūs geworht hæfde.
% mdparagraph:P259:end

\end{minipage}\par\addvspace{4pt}

\reserve{7\baselineskip}
% mdparagraph:P260:begin
\subsection*{A10}
% mdparagraph:P260:end

\par\begin{minipage}{\linewidth}\raggedright
% item: NEW-U12-A10
% mdparagraph:P261:begin
Translate as a report of intention. Which action is intended, and does the sentence assert that the intended action took place?
% mdparagraph:P261:end

% mdparagraph:P262:begin
Hē sægde ðæt hē wolde cuman.
% mdparagraph:P262:end

\end{minipage}\par\addvspace{4pt}

\reserve{7\baselineskip}
% mdparagraph:P263:begin
\subsection*{A11}
% mdparagraph:P263:end

\par\begin{minipage}{\linewidth}\raggedright
% item: NEW-U12-A11
% mdparagraph:P264:begin
The visit is still expected. Translate the first sentence in that setting. Translate the second as a statement of desire, and explain why English will does not always call for willan.
% mdparagraph:P264:end

% mdparagraph:P265:begin
Hē cymð.
% mdparagraph:P265:end

% mdparagraph:P266:begin
Hē wille cuman.
% mdparagraph:P266:end

\end{minipage}\par\addvspace{4pt}

\reserve{7\baselineskip}
% mdparagraph:P267:begin
\subsection*{C01}
% mdparagraph:P267:end

\par\begin{minipage}{\linewidth}\raggedright
% item: NEW-U12-C01
% mdparagraph:P268:begin
Translate: We have built the house. Use present habban with the selected past participle.
% mdparagraph:P268:end

\end{minipage}\par\addvspace{4pt}

\reserve{7\baselineskip}
% mdparagraph:P269:begin
\subsection*{C02}
% mdparagraph:P269:end

\par\begin{minipage}{\linewidth}\raggedright
% item: NEW-U12-C02
% mdparagraph:P270:begin
Translate: They had built the house. The building is earlier than another past event; use preterite habban with the selected past participle.
% mdparagraph:P270:end

\end{minipage}\par\addvspace{4pt}

\reserve{7\baselineskip}
% mdparagraph:P271:begin
\subsection*{C03}
% mdparagraph:P271:end

\par\begin{minipage}{\linewidth}\raggedright
% item: NEW-U12-C03
% mdparagraph:P272:begin
Translate: The men bound the horse. Then express the horse’s binding as an agentless passive, using preterite weorðan.
% mdparagraph:P272:end

\end{minipage}\par\addvspace{4pt}

\reserve{7\baselineskip}
% mdparagraph:P273:begin
\subsection*{C04}
% mdparagraph:P273:end

\par\begin{minipage}{\linewidth}\raggedright
% item: NEW-U12-C04
% mdparagraph:P274:begin
Translate: The monk says that he is here. Use the subjunctive to present his report. Then put the saying in the past, keeping his presence simultaneous with that past saying.
% mdparagraph:P274:end

\end{minipage}\par\addvspace{4pt}

\reserve{7\baselineskip}
% mdparagraph:P275:begin
\subsection*{C05}
% mdparagraph:P275:end

\par\begin{minipage}{\linewidth}\raggedright
% item: NEW-U12-C05
% mdparagraph:P276:begin
Translate: He said that he had built the house. Use the reporting subjunctive and preterite habban with the selected past participle.
% mdparagraph:P276:end

\end{minipage}\par\addvspace{4pt}

\reserve{7\baselineskip}
% mdparagraph:P277:begin
\subsection*{C06}
% mdparagraph:P277:end

\par\begin{minipage}{\linewidth}\raggedright
% item: NEW-U12-C06
% mdparagraph:P278:begin
Translate: He said that he intended to come. Use willan with the infinitive; do not assert that he actually came.
% mdparagraph:P278:end

\end{minipage}\par\addvspace{4pt}

\reserve{7\baselineskip}
% mdparagraph:P279:begin
\subsection*{C07}
% mdparagraph:P279:end

\par\begin{minipage}{\linewidth}\raggedright
% item: NEW-U12-C07
% mdparagraph:P280:begin
Address one person: Prepare the food. Then translate the reply: I have prepared the food. Use gierwan in both.
% mdparagraph:P280:end

\end{minipage}\par\addvspace{4pt}

\reserve{7\baselineskip}
% mdparagraph:P281:begin
\subsection*{C08}
% mdparagraph:P281:end

\par\begin{minipage}{\linewidth}\raggedright
% item: NEW-U12-C08
% mdparagraph:P282:begin
Translate: The bridge was built. Then a man came. Use an agentless passive with wesan for the completed bridge, followed by a simple preterite for the arrival.
% mdparagraph:P282:end

\end{minipage}\par\addvspace{4pt}

\reserve{7\baselineskip}
% mdparagraph:P283:begin
\section*{10. Cook, Group VI: Exercises 19–21}
% mdparagraph:P283:end

% teaching: NEW-U12-LOCAL-HELPS
% mdparagraph:P284:begin
Translate all the following into Old English. The English wording, parentheses, omissions and section references are Cook’s. Directions and supplied help are separate. \textbf{Om.} means omit the indicated English word in the Old English construction; it does not authorize dropping the meaning of the sentence. \textbf{Sing.} specifies a singular addressee; \textbf{pret.} requires a preterite. Cook’s numbered references are to his grammar, not the sections of this unit.
% mdparagraph:P284:end

% mdparagraph:P285:begin
Use the taught vocabulary and the stated constructions. A valid alternative may differ from a selected target; when a direction specifies a construction or addressee, satisfy that direction. The fixed expressions and exceptional local forms printed as help remain available during the exercise.
% mdparagraph:P285:end

\reserve{7\baselineskip}
% mdparagraph:P286:begin
\subsection*{Exercise 19 — printed pp. 9–10 /\allowbreak{} PDF pp. 17–18}
% mdparagraph:P286:end

\reserve{7\baselineskip}
% mdparagraph:P287:begin
\subsection*{19.1}
% mdparagraph:P287:end

\par\begin{minipage}{\linewidth}\raggedright
% item: EX-VI-19-01
% mdparagraph:P288:begin
At length the king ceased to bear rule.
% mdparagraph:P288:end

% mdparagraph:P289:begin
Direction: Use geswīcan with a genitive for ceasing from rule.
% mdparagraph:P289:end

\end{minipage}\par\addvspace{4pt}

% teaching: NEW-U12-REUSED-U12-COOK-HELP-19-01
% teaching: NEW-U12-GRAM-U12-COOK-CONSTRUCTION-19-01
% mdparagraph:P290:begin
\textbf{Supplied help.} Use the fixed phrase æt nīehstan, at length. Its ending is supplied; no general superlative system is required.
% mdparagraph:P290:end

\reserve{7\baselineskip}
% mdparagraph:P291:begin
\subsection*{19.2}
% mdparagraph:P291:end

\par\begin{minipage}{\linewidth}\raggedright
% item: EX-VI-19-02
% mdparagraph:P292:begin
He had the kingdom forty years.
% mdparagraph:P292:end

% mdparagraph:P293:begin
Direction: Distinguish possessing the kingdom from the new perfect construction.
% mdparagraph:P293:end

\end{minipage}\par\addvspace{4pt}

% teaching: NEW-U12-GRAM-U12-COOK-CONSTRUCTION-19-02
% mdparagraph:P294:begin
\textbf{Supplied help.} Use fēowertig gēara, forty years. This whole duration phrase is supplied; the higher numeral takes the genitive gēara. The duration construction is discussed in this unit.
% mdparagraph:P294:end

\reserve{7\baselineskip}
% mdparagraph:P295:begin
\subsection*{19.3}
% mdparagraph:P295:end

\par\begin{minipage}{\linewidth}\raggedright
% item: EX-VI-19-03
% mdparagraph:P296:begin
When were the bridges built?
% mdparagraph:P296:end

% mdparagraph:P297:begin
Direction: Use a past agentless passive, with the bridges as subject.
% mdparagraph:P297:end

\end{minipage}\par\addvspace{4pt}

\reserve{7\baselineskip}
% mdparagraph:P298:begin
\subsection*{19.4}
% mdparagraph:P298:end

\par\begin{minipage}{\linewidth}\raggedright
% item: EX-VI-19-04
% mdparagraph:P299:begin
To-day thou shalt be with me.
% mdparagraph:P299:end

% mdparagraph:P300:begin
Direction: Use a taught form of be in its future context.
% mdparagraph:P300:end

\end{minipage}\par\addvspace{4pt}

\reserve{7\baselineskip}
% mdparagraph:P301:begin
\subsection*{19.5}
% mdparagraph:P301:end

\par\begin{minipage}{\linewidth}\raggedright
% item: EX-VI-19-05
% mdparagraph:P302:begin
What (om. did) the Romans build (pret.) in Britain?
% mdparagraph:P302:end

% mdparagraph:P303:begin
Direction: Keep the explicit preterite requirement.
% mdparagraph:P303:end

\end{minipage}\par\addvspace{4pt}

% teaching: NEW-U12-GRAM-U12-COOK-CONSTRUCTION-19-05
% mdparagraph:P304:begin
\textbf{Supplied help.} Rōmane, Romans, is the supplied nominative plural; on Breotone, in Britain, supplies the place-name dative. Retain the preterite target.
% mdparagraph:P304:end

\reserve{7\baselineskip}
% mdparagraph:P305:begin
\subsection*{19.6}
% mdparagraph:P305:end

\par\begin{minipage}{\linewidth}\raggedright
% item: EX-VI-19-06
% mdparagraph:P306:begin
I send my angel before (beforan, 166) thy face.
% mdparagraph:P306:end

% mdparagraph:P307:begin
Direction: Preserve the original face wording and the specified beforan phrase.
% mdparagraph:P307:end

\end{minipage}\par\addvspace{4pt}

\reserve{7\baselineskip}
% mdparagraph:P308:begin
\subsection*{19.7}
% mdparagraph:P308:end

\par\begin{minipage}{\linewidth}\raggedright
% item: EX-VI-19-07
% mdparagraph:P309:begin
Come (sing.), and bring those letters that you wrote.
% mdparagraph:P309:end

% mdparagraph:P310:begin
Direction: Treat you in the relative clause as the single person addressed by the commands.
% mdparagraph:P310:end

\end{minipage}\par\addvspace{4pt}

\reserve{7\baselineskip}
% mdparagraph:P311:begin
\subsection*{19.8}
% mdparagraph:P311:end

\par\begin{minipage}{\linewidth}\raggedright
% item: EX-VI-19-08
% mdparagraph:P312:begin
I promise you obedience, if you help (196. d) me.
% mdparagraph:P312:end

% mdparagraph:P313:begin
Direction: Cook §196.d calls for the selected subjunctive in the if-clause. Keep singular or plural you consistent.
% mdparagraph:P313:end

\end{minipage}\par\addvspace{4pt}

\reserve{7\baselineskip}
% mdparagraph:P314:begin
\subsection*{19.9}
% mdparagraph:P314:end

\par\begin{minipage}{\linewidth}\raggedright
% item: EX-VI-19-09
% mdparagraph:P315:begin
Send me (om. a) large army.
% mdparagraph:P315:end

% mdparagraph:P316:begin
Direction: Address one person; distinguish the thing sent from its recipient.
% mdparagraph:P316:end

\end{minipage}\par\addvspace{4pt}

\reserve{7\baselineskip}
% mdparagraph:P317:begin
\subsection*{19.10}
% mdparagraph:P317:end

\par\begin{minipage}{\linewidth}\raggedright
% item: EX-VI-19-10
% mdparagraph:P318:begin
I fight with my enemies.
% mdparagraph:P318:end

% mdparagraph:P319:begin
Direction: Retain the earlier selected feohtan wið + accusative opponent construction.
% mdparagraph:P319:end

\end{minipage}\par\addvspace{4pt}

\reserve{7\baselineskip}
% mdparagraph:P320:begin
\subsection*{19.11}
% mdparagraph:P320:end

\par\begin{minipage}{\linewidth}\raggedright
% item: EX-VI-19-11
% mdparagraph:P321:begin
Expel your foes.
% mdparagraph:P321:end

% mdparagraph:P322:begin
Direction: Address more than one person and make the possessive agree with the foes.
% mdparagraph:P322:end

\end{minipage}\par\addvspace{4pt}

\reserve{7\baselineskip}
% mdparagraph:P323:begin
\subsection*{19.12}
% mdparagraph:P323:end

\par\begin{minipage}{\linewidth}\raggedright
% item: EX-VI-19-12
% mdparagraph:P324:begin
The soldiers went home.
% mdparagraph:P324:end

% mdparagraph:P325:begin
Direction: Use the directional adverb home without adding a preposition.
% mdparagraph:P325:end

\end{minipage}\par\addvspace{4pt}

\reserve{7\baselineskip}
% mdparagraph:P326:begin
\subsection*{19.13}
% mdparagraph:P326:end

\par\begin{minipage}{\linewidth}\raggedright
% item: EX-VI-19-13
% mdparagraph:P327:begin
He will come with (om. a) fleet.
% mdparagraph:P327:end

% mdparagraph:P328:begin
Direction: Express future reference without treating willan as compulsory.
% mdparagraph:P328:end

\end{minipage}\par\addvspace{4pt}

\reserve{7\baselineskip}
% mdparagraph:P329:begin
\subsection*{19.14}
% mdparagraph:P329:end

\par\begin{minipage}{\linewidth}\raggedright
% item: EX-VI-19-14
% mdparagraph:P330:begin
Slay (sing.) all that you find (gemētan).
% mdparagraph:P330:end

% mdparagraph:P331:begin
Direction: Address one person. For this target use all those who with a relative clause.
% mdparagraph:P331:end

\end{minipage}\par\addvspace{4pt}

\reserve{7\baselineskip}
% mdparagraph:P332:begin
\subsection*{19.15}
% mdparagraph:P332:end

\par\begin{minipage}{\linewidth}\raggedright
% item: EX-VI-19-15
% mdparagraph:P333:begin
The harvest is rich in fruits.
% mdparagraph:P333:end

% mdparagraph:P334:begin
Direction: Use the rich-in construction taught in §8.
% mdparagraph:P334:end

\end{minipage}\par\addvspace{4pt}

\reserve{7\baselineskip}
% mdparagraph:P335:begin
\subsection*{19.16}
% mdparagraph:P335:end

\par\begin{minipage}{\linewidth}\raggedright
% item: EX-VI-19-16
% mdparagraph:P336:begin
Seedtime and harvest never shall fail (geswīcan).
% mdparagraph:P336:end

% mdparagraph:P337:begin
Direction: Use objectless geswīcan, the plural subject, and the taught never construction.
% mdparagraph:P337:end

\end{minipage}\par\addvspace{4pt}

\reserve{7\baselineskip}
% mdparagraph:P338:begin
\subsection*{19.17}
% mdparagraph:P338:end

\par\begin{minipage}{\linewidth}\raggedright
% item: EX-VI-19-17
% mdparagraph:P339:begin
The year has four seasons.
% mdparagraph:P339:end

% mdparagraph:P340:begin
Direction: Use ordinary possession, not an auxiliary perfect.
% mdparagraph:P340:end

\end{minipage}\par\addvspace{4pt}

\reserve{7\baselineskip}
% mdparagraph:P341:begin
\subsection*{Exercise 20 — printed p. 10 /\allowbreak{} PDF p. 18}
% mdparagraph:P341:end

\reserve{7\baselineskip}
% mdparagraph:P342:begin
\subsection*{20.1}
% mdparagraph:P342:end

\par\begin{minipage}{\linewidth}\raggedright
% item: EX-VI-20-01
% mdparagraph:P343:begin
The earth is (om. the) Lord’s.
% mdparagraph:P343:end

% mdparagraph:P344:begin
Direction: Use a genitive possessor after be.
% mdparagraph:P344:end

\end{minipage}\par\addvspace{4pt}

\reserve{7\baselineskip}
% mdparagraph:P345:begin
\subsection*{20.2}
% mdparagraph:P345:end

\par\begin{minipage}{\linewidth}\raggedright
% item: EX-VI-20-02
% mdparagraph:P346:begin
The sea saw, and it fled.
% mdparagraph:P346:end

% mdparagraph:P347:begin
Direction: Retain the selected masculine gender of sǣ and the objectless seeing sense.
% mdparagraph:P347:end

\end{minipage}\par\addvspace{4pt}

\reserve{7\baselineskip}
% mdparagraph:P348:begin
\subsection*{20.3}
% mdparagraph:P348:end

\par\begin{minipage}{\linewidth}\raggedright
% item: EX-VI-20-03
% mdparagraph:P349:begin
He shall flee the iron weapon.
% mdparagraph:P349:end

% mdparagraph:P350:begin
Direction: The thing avoided is accusative; distinguish the adjective iron from the noun.
% mdparagraph:P350:end

\end{minipage}\par\addvspace{4pt}

\reserve{7\baselineskip}
% mdparagraph:P351:begin
\subsection*{20.4}
% mdparagraph:P351:end

\par\begin{minipage}{\linewidth}\raggedright
% item: EX-VI-20-04
% mdparagraph:P352:begin
He that (sē ðe) is not (om. the) shepherd, and owns not (127) the sheep, when he sees the wolf, then flees he, and abandons the sheep, and the wolf seizes the sheep.
% mdparagraph:P352:end

% mdparagraph:P353:begin
Direction: Translate the whole relative–temporal–main-clause sequence and preserve its changing subjects.
% mdparagraph:P353:end

\end{minipage}\par\addvspace{4pt}

% teaching: NEW-U12-GRAM-U12-COOK-CONSTRUCTION-20-04
% teaching: NEW-U12-GRAM-u12-negative-agan-local
% mdparagraph:P354:begin
\textbf{Supplied help.} Use nāh, does not own, for the negative possession clause. Cook’s hint sē ðe is retained as printed; the model may use the earlier selected Se ðe, he who. The long-vowel sē is a supplied source spelling, not a newly required paradigm.
% mdparagraph:P354:end

% teaching: NEW-U12-REUSED-U12-COOK-HELP-NIS
% teaching: NEW-U12-GRAM-U12-COOK-CONSTRUCTION-NIS
% mdparagraph:P355:begin
\textbf{Supplied help.} Reuse the supplied contraction nis, is not, for the negative be-clause. This local form does not teach the full negative verb system.
% mdparagraph:P355:end

\reserve{7\baselineskip}
% mdparagraph:P356:begin
\subsection*{20.5}
% mdparagraph:P356:end

\par\begin{minipage}{\linewidth}\raggedright
% item: EX-VI-20-05
% mdparagraph:P357:begin
The day is thine, the night also is thine; thou hast prepared the sun and the moon; thou hast set the borders of the earth; thou hast made summer and winter.
% mdparagraph:P357:end

% mdparagraph:P358:begin
Direction: Use agreeing possessive predicates, then habban with the selected past participle for each completed action.
% mdparagraph:P358:end

\end{minipage}\par\addvspace{4pt}

\reserve{7\baselineskip}
% mdparagraph:P359:begin
\subsection*{20.6}
% mdparagraph:P359:end

\par\begin{minipage}{\linewidth}\raggedright
% item: EX-VI-20-06
% mdparagraph:P360:begin
Let us build these cities, and make about (ābūtan, 172) them walls, and towers, and gates, while the land is yet (gīet) before us.
% mdparagraph:P360:end

% mdparagraph:P361:begin
Direction: Use first-plural subjunctive exhortations with gewyrcean for both building clauses, the specified ābūtan phrase, and the factual while-clause.
% mdparagraph:P361:end

\end{minipage}\par\addvspace{4pt}

\reserve{7\baselineskip}
% mdparagraph:P362:begin
\subsection*{20.7}
% mdparagraph:P362:end

\par\begin{minipage}{\linewidth}\raggedright
% item: EX-VI-20-07
% mdparagraph:P363:begin
These nations are more than I; how can I expel them?
% mdparagraph:P363:end

% mdparagraph:P364:begin
Direction: Use indeclinable quantitative mā and magan with the plain infinitive of expel.
% mdparagraph:P364:end

\end{minipage}\par\addvspace{4pt}

\reserve{7\baselineskip}
% mdparagraph:P365:begin
\subsection*{20.8}
% mdparagraph:P365:end

\par\begin{minipage}{\linewidth}\raggedright
% item: EX-VI-20-08
% mdparagraph:P366:begin
Go to the sea, and cast a fishhook, and take the first (ǣrest) fish, and open (geopenian) his mouth (mūð, 43).
% mdparagraph:P366:end

% mdparagraph:P367:begin
Direction: Keep all the commands singular and preserve the fish’s masculine reference.
% mdparagraph:P367:end

\end{minipage}\par\addvspace{4pt}

% teaching: NEW-U12-GRAM-U12-COOK-CONSTRUCTION-20-08
% mdparagraph:P368:begin
\textbf{Supplied help.} Cook prints first (ǣrest). Treat this as a lexical hint: use the previously taught ordinal ǣresta and the needed agreeing form ǣrestan in ðone ǣrestan fisc. The uninflected printed hint is preserved in the prompt.
% mdparagraph:P368:end

\reserve{7\baselineskip}
% mdparagraph:P369:begin
\subsection*{20.9}
% mdparagraph:P369:end

\par\begin{minipage}{\linewidth}\raggedright
% item: EX-VI-20-09
% mdparagraph:P370:begin
Thou wilt find a penny; take that, and give (sęllan) for me and for thee.
% mdparagraph:P370:end

% mdparagraph:P371:begin
Direction: Keep the penny masculine; use for + dative in both beneficiary phrases.
% mdparagraph:P371:end

\end{minipage}\par\addvspace{4pt}

\reserve{7\baselineskip}
% mdparagraph:P372:begin
\subsection*{Exercise 21 — printed pp. 10–11 /\allowbreak{} PDF pp. 18–19}
% mdparagraph:P372:end

\reserve{7\baselineskip}
% mdparagraph:P373:begin
\subsection*{21.1}
% mdparagraph:P373:end

\par\begin{minipage}{\linewidth}\raggedright
% item: EX-VI-21-01
% mdparagraph:P374:begin
Then said he, She is not dead.
% mdparagraph:P374:end

% mdparagraph:P375:begin
Direction: Retain the past report and the present state in direct speech.
% mdparagraph:P375:end

% mdparagraph:P376:begin
\textbf{Supplied help.} Reuse the supplied contraction nis, is not, for the negative be-clause. This local form does not teach the full negative verb system.
% mdparagraph:P376:end

\end{minipage}\par\addvspace{4pt}

\reserve{7\baselineskip}
% mdparagraph:P377:begin
\subsection*{21.2}
% mdparagraph:P377:end

\par\begin{minipage}{\linewidth}\raggedright
% item: EX-VI-21-02
% mdparagraph:P378:begin
Then took he her hand, and said, Maiden, arise.
% mdparagraph:P378:end

% mdparagraph:P379:begin
Direction: Preserve the feminine possessor and address the maiden with the singular command of rīsan.
% mdparagraph:P379:end

\end{minipage}\par\addvspace{4pt}

\reserve{7\baselineskip}
% mdparagraph:P380:begin
\subsection*{21.3}
% mdparagraph:P380:end

\par\begin{minipage}{\linewidth}\raggedright
% item: EX-VI-21-03
% mdparagraph:P381:begin
The stone wall is seven feet broad, and twenty feet high.
% mdparagraph:P381:end

% mdparagraph:P382:begin
Direction: Use both genitive measure phrases and adjectives agreeing with the wall.
% mdparagraph:P382:end

\end{minipage}\par\addvspace{4pt}

% teaching: NEW-U12-GRAM-U12-COOK-CONSTRUCTION-21-03
% mdparagraph:P383:begin
\textbf{Supplied help.} Twēntig means twenty. Use the supplied measure phrase twēntig fōta; the newly taught genitive of extent also governs seofon fōta.
% mdparagraph:P383:end

\reserve{7\baselineskip}
% mdparagraph:P384:begin
\subsection*{21.4}
% mdparagraph:P384:end

\par\begin{minipage}{\linewidth}\raggedright
% item: EX-VI-21-04
% mdparagraph:P385:begin
I took leave of her (acc.), and announced (om. to) her that I never more would visit her.
% mdparagraph:P385:end

% mdparagraph:P386:begin
Direction: Use grētan, cȳðan and an intended later genēosian with willan. Use the reporting subjunctive; do not assert that a visit occurred.
% mdparagraph:P386:end

\end{minipage}\par\addvspace{4pt}

% teaching: NEW-U12-GRAM-U12-COOK-CONSTRUCTION-21-04
% mdparagraph:P387:begin
\textbf{Supplied help.} Use the supplied nolde, the contracted negative preterite subjunctive singular of willan, in the reported later intention. Nǣfre mā adds never again; it combines with the negative finite verb. This one supplied contraction does not assign the entire negative paradigm.
% mdparagraph:P387:end

\reserve{7\baselineskip}
% mdparagraph:P388:begin
\subsection*{21.5}
% mdparagraph:P388:end

\par\begin{minipage}{\linewidth}\raggedright
% item: EX-VI-21-05
% mdparagraph:P389:begin
The righteous (om. ones) are as bold as (om. a) lion.
% mdparagraph:P389:end

% mdparagraph:P390:begin
Direction: Use the weak adjective with an article for the righteous people, and separated swā … adjective … swā.
% mdparagraph:P390:end

\end{minipage}\par\addvspace{4pt}

% teaching: NEW-U12-GRAM-U12-COOK-CONSTRUCTION-21-05
% mdparagraph:P391:begin
\textbf{Supplied help.} Use the supplied nominative lēo, lion, after the second swā. No full paradigm of this loanword is assigned.
% mdparagraph:P391:end

\reserve{7\baselineskip}
% mdparagraph:P392:begin
\subsection*{21.6}
% mdparagraph:P392:end

\par\begin{minipage}{\linewidth}\raggedright
% item: EX-VI-21-06
% mdparagraph:P393:begin
The Britons were cowardly.
% mdparagraph:P393:end

% mdparagraph:P394:begin
Direction: Make the predicate adjective plural.
% mdparagraph:P394:end

\end{minipage}\par\addvspace{4pt}

% teaching: NEW-U12-GRAM-U12-COOK-CONSTRUCTION-21-06
% mdparagraph:P395:begin
\textbf{Supplied help.} Bryttas, Britons, is supplied as a nominative plural.
% mdparagraph:P395:end

\reserve{7\baselineskip}
% mdparagraph:P396:begin
\subsection*{21.7}
% mdparagraph:P396:end

\par\begin{minipage}{\linewidth}\raggedright
% item: EX-VI-21-07
% mdparagraph:P397:begin
Come hither, fair maiden.
% mdparagraph:P397:end

% mdparagraph:P398:begin
Direction: Use the weak adjective in singular direct address.
% mdparagraph:P398:end

\end{minipage}\par\addvspace{4pt}

\reserve{7\baselineskip}
% mdparagraph:P399:begin
\subsection*{21.8}
% mdparagraph:P399:end

\par\begin{minipage}{\linewidth}\raggedright
% item: EX-VI-21-08
% mdparagraph:P400:begin
Happy is that people whose God is the Lord.
% mdparagraph:P400:end

% mdparagraph:P401:begin
Direction: Distinguish the case of the antecedent people from the genitive function of whose.
% mdparagraph:P401:end

\end{minipage}\par\addvspace{4pt}

\reserve{7\baselineskip}
% mdparagraph:P402:begin
\subsection*{21.9}
% mdparagraph:P402:end

\par\begin{minipage}{\linewidth}\raggedright
% item: EX-VI-21-09
% mdparagraph:P403:begin
Learn of (æt) me, for I am merciful and lowly (ēaðmōd) in heart, and ye shall find (gemētan) rest (ręst, 51. b) (om. unto) your souls.
% mdparagraph:P403:end

% mdparagraph:P404:begin
Direction: Address more than one person. Use the supplied original æt hint, on + dative for respect, and dative beneficiaries after finding rest.
% mdparagraph:P404:end

\end{minipage}\par\addvspace{4pt}

\reserve{7\baselineskip}
% mdparagraph:P405:begin
\subsection*{21.10}
% mdparagraph:P405:end

\par\begin{minipage}{\linewidth}\raggedright
% item: EX-VI-21-10
% mdparagraph:P406:begin
Now I send you as sheep among (gemang, 172) wolves.
% mdparagraph:P406:end

% mdparagraph:P407:begin
Direction: Use the single-swā nominal comparison and the selected gemang + accusative phrase.
% mdparagraph:P407:end

\end{minipage}\par\addvspace{4pt}

\reserve{7\baselineskip}
% mdparagraph:P408:begin
\subsection*{21.11}
% mdparagraph:P408:end

\par\begin{minipage}{\linewidth}\raggedright
% item: EX-VI-21-11
% mdparagraph:P409:begin
When they pursue (ēhtan) you in this city, flee into (on) (om. an)other.
% mdparagraph:P409:end

% mdparagraph:P410:begin
Direction: Use the selected accusative person pursued, dative location, and plural flight command.
% mdparagraph:P410:end

\end{minipage}\par\addvspace{4pt}

% teaching: NEW-U12-GRAM-U12-COOK-CONSTRUCTION-21-11
% teaching: NEW-U12-GRAM-u12-on-elliptic-city
% mdparagraph:P411:begin
\textbf{Supplied help.} Use on ōðre, into another, with feminine accusative singular ōðre and an understood ceastre, city. Cook prints (om. an)other: omit the article an, not the sense another. The city omission is supplied locally.
% mdparagraph:P411:end

\reserve{7\baselineskip}
% mdparagraph:P412:begin
\subsection*{21.12}
% mdparagraph:P412:end

\par\begin{minipage}{\linewidth}\raggedright
% item: EX-VI-21-12
% mdparagraph:P413:begin
I take leave of my native land.
% mdparagraph:P413:end

% mdparagraph:P414:begin
Direction: Use present grētan with the accusative native land.
% mdparagraph:P414:end

\end{minipage}\par\addvspace{4pt}

\reserve{7\baselineskip}
% mdparagraph:P415:begin
\section*{11. A house and a return}
% mdparagraph:P415:end

% mdparagraph:P416:begin
This connected account is new teaching material. Its narrative and direct speech use different viewpoints; the report of an intention does not by itself assert that the intended action happened.
% mdparagraph:P416:end

% composition: NEW-U12-TEXT-A
% mdparagraph:P417:begin
\par\begin{minipage}{\linewidth}\raggedright
\readingtext{Se mann hæfde ān hūs geworht. Ðā cōm his sunu, and se sunu geseah ðæt hūs. Se fæder cwæð: “Nū hæbbe ic ðæt hūs geworht.” Se sunu sægde ðæt hē wolde faran. Ðā se sunu gewāt, ðā sæt se fæder on ðǣm hūse. Ðā cōm se sunu, and hē sealde his fæder męte.}
\end{minipage}\par
% mdparagraph:P417:end

\reserve{7\baselineskip}
% mdparagraph:P418:begin
\subsection*{R01}
% mdparagraph:P418:end

\par\begin{minipage}{\linewidth}\raggedright
% item: NEW-U12-R01
% mdparagraph:P419:begin
Translate the whole new account, A house and a return. Preserve its changes between narrative, direct speech and reported intention.
% mdparagraph:P419:end

\end{minipage}\par\addvspace{4pt}

\reserve{7\baselineskip}
% mdparagraph:P420:begin
\subsection*{R02}
% mdparagraph:P420:end

\par\begin{minipage}{\linewidth}\raggedright
% item: NEW-U12-R02
% mdparagraph:P421:begin
Put these events in order: the son’s first arrival; building the house; the son’s departure; the father’s announcement; the son’s return with food. Which verb phrase places one event before the story’s first arrival?
% mdparagraph:P421:end

\end{minipage}\par\addvspace{4pt}

\reserve{7\baselineskip}
% mdparagraph:P422:begin
\subsection*{R03}
% mdparagraph:P422:end

\par\begin{minipage}{\linewidth}\raggedright
% item: NEW-U12-R03
% mdparagraph:P423:begin
Identify the referents of ic in the father’s speech, hē in the reported intention, and hē in the last sentence. Is the journey an intention only, or does the following narrative show a departure? Cite the relevant words.
% mdparagraph:P423:end

\end{minipage}\par\addvspace{4pt}

\reserve{7\baselineskip}
% mdparagraph:P424:begin
\subsection*{R04}
% mdparagraph:P424:end

\par\begin{minipage}{\linewidth}\raggedright
% item: NEW-U12-R04
% mdparagraph:P425:begin
In Nū hæbbe ic ðæt hūs geworht, identify the subject and object. Express the house’s completed state in an agentless sentence with present wesan and the same participle.
% mdparagraph:P425:end

\end{minipage}\par\addvspace{4pt}

\reserve{7\baselineskip}
% mdparagraph:P426:begin
\section*{12. First Roman relief and the renewed invasion}
% mdparagraph:P426:end

% mdparagraph:P427:begin
This exact 155-word extract comes from Cook’s \emph{First Book}, Reader VI, \textbf{The Invasion of Britain by the Picts and Scots}, printed pp. 145–146 /\allowbreak{} PDF pp. 161–162. It begins at editorial prose line 1 of p. 145, includes all seven prose lines there, and continues through foriermdon at the end of the first sentence on editorial prose line 11 of p. 146. These editorial counts differ from Cook’s printed marginal line numbers. The following English context is separate from the Old English.
% mdparagraph:P427:end

% mdparagraph:P428:begin
The preceding passage says that Roman rule in Britain had ended and describes the Britons within the old wall. In this extract the Britons seek military help from Rome after raids by the Picts and Scots. This context is an editorial English summary, not translated words added to the Old English excerpt.
% mdparagraph:P428:end

% teaching: NEW-U12-SOURCE-HELPS
\par\begin{minipage}{\linewidth}\raggedright
% source: FB-VI-PICTS-SCOTS-RELIEF-01
% mdparagraph:P429:begin
\readingtext{Ðā ongunnon twā ðēoda, Pihtas norðan, and Scottas westan, hīe onwinnan, and hiera ǣhta niman and hęrgian; and hīe fela gēara iermdon and hīendon. Ðā, on ðǣre unstilnesse, onsęndon hīe ǣrendwrecan tō Rome mid gewritum and wēpendre bēne; him fultumes bǣdon, and him gehēton ēaðmōde hīernesse and singāle underðēodnesse, gif hīe him gefultumoden ðæt hīe meahten hiera fīend oferwinnan. Ðā onsęndon hīe him micelne hęre tō fultume, and, sōna ðæs-ðe hīe on ðis īgland cōmon, ðā campedon hīe wið hiera fēondum, and him micel wæl ongeslōgon, and of hiera gemǣrum ādrifon and āflīemdon; and lǣrdon ðæt hīe fæsten worhten him tō gebeorge wið hiera fēondum; and swā, mid micle sige, hām fōron.}
% mdparagraph:P429:end

% mdparagraph:P430:begin
\readingtext{Ðā ðæt ðā ongēaton ðā ǣrran gewinnan, ðæt se Rōmanisca hęre wæs onweg gewiten, ðā cōmon hīe sōna mid sciphęre on hiera landgemǣru, and slōgon ealle and cwealdon ðæt hīe gemētton, and swā-swā rīpe ierðe fortrǣdon and fornōmon, and hīe mid ealle foriermdon.}
% mdparagraph:P430:end

\end{minipage}\par
\reserve{7\baselineskip}
% mdparagraph:P431:begin
\subsection*{Reading help}
% mdparagraph:P431:end

% mdparagraph:P432:begin
Use all the following help while translating and analyzing the passage. Extra historical vocabulary and the expressly supplied constructions are not independent-recall material. Earlier forms retain their ordinary uses; new senses are identified where needed. The note numbers cited below refer to Cook’s original apparatus in §13.
% mdparagraph:P432:end

% teaching: NEW-U12-GRAM-NEW-U12-SRC-01
% support: U12-SRC-HELP-01
% mdparagraph:P433:begin
\textbf{1. Ðā ongunnon ... hęrgian.} Ongunnon, began, takes the three infinitives onwinnan, attack, niman, take, and hęrgian, plunder. Pihtas are the Picts, Scottas the Scots; norðan and westan mean from the north and from the west. Hīe refers to the Britons introduced in the separate context; ǣhta means possessions.
% mdparagraph:P433:end

% teaching: NEW-U12-GRAM-NEW-U12-SRC-02
% support: U12-SRC-HELP-02
% mdparagraph:P434:begin
\textbf{2. and hīe fela gēara iermdon and hīendon.} Iermdon and hīendon mean afflicted and oppressed. Fela gēara means for many years, with genitive after fela. Hīe can be the Britons as object with the raiders continuing as subject, or the raiders as subject with a Briton object understood; the events do not change. Do not infer a unique case from the spelling.
% mdparagraph:P434:end

% teaching: NEW-U12-GRAM-NEW-U12-SRC-03
% support: U12-SRC-HELP-03
% mdparagraph:P435:begin
\textbf{3. Ðā, on ðǣre unstilnesse ... bēne.} On ðǣre unstilnesse, amid that unrest. Onsęndon is the exact prefixed sent; the Britons are now the subject. Ǣrendwrecan, messengers; tō Rome, to Rome. Mid gewritum and wēpendre bēne means with letters and a tearful petition. Wēpendre is a present participle, strong dative feminine singular agreeing with bēne.
% mdparagraph:P435:end

% teaching: NEW-U12-GRAM-NEW-U12-SRC-04
% support: U12-SRC-HELP-04
% mdparagraph:P436:begin
\textbf{4. him fultumes bǣdon ... underðēodnesse.} In him fultumes bǣdon, him names the Romans asked and fultumes is the genitive thing requested: help. Hīernes is obedience; singāl is continual; underðēodnes is subjection. Ēaðmōde and singāle are accusative feminine singular. The promise is made by the Britons to the Romans.
% mdparagraph:P436:end

% teaching: NEW-U12-GRAM-NEW-U12-SRC-05
% teaching: NEW-U12-GRAM-u12-source-reported-condition
% support: U12-SRC-HELP-05
% mdparagraph:P437:begin
\textbf{5. gif hīe him gefultumoden ... oferwinnan.} Gefultumoden is a preterite subjunctive: if they, the Romans, would help them, the Britons. The following ðæt clause gives the purpose, so that they might overcome their enemies. Meahten is the known preterite subjunctive plural of magan; oferwinnan means overcome. This whole reported condition/\allowbreak{}purpose sequence is supported locally.
% mdparagraph:P437:end

% teaching: NEW-U12-GRAM-NEW-U12-SRC-06
% support: U12-SRC-HELP-06
% mdparagraph:P438:begin
\textbf{6. Ðā onsęndon ... tō fultume.} Here the subject of onsęndon is the Romans, while him refers to the Britons. Tō fultume means as help/\allowbreak{}to their aid. The Old English says they sent an army in the active voice; Cook compares a passive Latin expression. Preserve the actual Old English voice when analyzing it.
% mdparagraph:P438:end

% teaching: NEW-U12-GRAM-NEW-U12-SRC-07
% teaching: NEW-U12-REUSED-U12-SRC-HELP-07
% support: U12-SRC-HELP-07
% mdparagraph:P439:begin
\textbf{7. and, sōna ðæs-ðe ... cōmon, ðā.} Sōna ðæs-ðe means as soon as. On ðis īgland has the accusative goal, to/\allowbreak{}onto this island. Hīe refers to the Roman force. The following ðā resumes the narrative.
% mdparagraph:P439:end

% teaching: NEW-U12-GRAM-NEW-U12-SRC-08
% teaching: NEW-U12-REUSED-U12-SRC-HELP-08
% support: U12-SRC-HELP-08
% mdparagraph:P440:begin
\textbf{8. campedon ... āflīemdon.} Campedon, fought; wið hiera fēondum, against their enemies, with the source dative after wið. Him micel wæl ongeslōgon means inflicted great slaughter on them. Of hiera gemǣrum means out of their borders, with dative after of; the borders belong to the Britons, Rome’s allies. Āflīemdon means put to flight, reinforcing ādrifon. The enemy object carries on through the coordinated verbs.
% mdparagraph:P440:end

% teaching: NEW-U12-GRAM-NEW-U12-SRC-09
% teaching: NEW-U12-REUSED-U12-SRC-HELP-09
% support: U12-SRC-HELP-09
% mdparagraph:P441:begin
\textbf{9. lǣrdon ðæt ... fēondum.} Lǣrdon means advised/\allowbreak{}instructed. Fæsten is a fortification; worhten is the known unprefixed wyrcean preterite subjunctive, here make/\allowbreak{}build in the advice clause. Him tō gebeorge means for their own protection. The Romans give the advice and the Britons are to build. The source’s wið + dative remains locally supplied.
% mdparagraph:P441:end

% teaching: NEW-U12-GRAM-NEW-U12-SRC-10
% support: U12-SRC-HELP-10
% mdparagraph:P442:begin
\textbf{10. and swā, mid micle sige, hām fōron.} Sige means victory. Mid micle sige is with great victory: the adjective has the instrumental micle beside the dative noun, as Cook’s note directs. Treat this as a supplied source construction. Hām is the directional home, without a preposition.
% mdparagraph:P442:end

% teaching: NEW-U12-GRAM-NEW-U12-SRC-11
% support: U12-SRC-HELP-11
% mdparagraph:P443:begin
\textbf{11. Ðā ðæt ðā ongēaton ðā ǣrran gewinnan, ðæt.} Read the three ðā forms in order as when, then, the. The first ðæt is an anticipatory object, this/\allowbreak{}that, expanded by the later that-clause. Ongēaton means realized; ǣrran gewinnan means former enemies. Gewinnan is a weak masculine plural noun, distinct from the previously supplied neuter gewinn, trouble.
% mdparagraph:P443:end

% teaching: NEW-U12-GRAM-NEW-U12-SRC-12
% support: U12-SRC-HELP-12
% mdparagraph:P444:begin
\textbf{12. se Rōmanisca hęre wæs onweg gewiten.} Rōmanisca is the weak masculine adjective Roman. Wæs onweg gewiten means had gone away: gewiten is the past participle of intransitive gewītan. This is an intransitive completed-state/\allowbreak{}perfect expression, not a passive. The source has plain-e onweg; retain it beside the earlier selected onwęg.
% mdparagraph:P444:end

% teaching: NEW-U12-GRAM-NEW-U12-SRC-13
% teaching: NEW-U12-REUSED-U12-SRC-HELP-13
% support: U12-SRC-HELP-13
% mdparagraph:P445:begin
\textbf{13. ðā cōmon ... landgemǣru.} The returning subject is the Picts and Scots. Sōna means immediately; mid sciphęre, with a fleet. On hiera landgemǣru has accusative plural for movement into their territory; the territory is the Britons’.
% mdparagraph:P445:end

% teaching: NEW-U12-GRAM-NEW-U12-SRC-14
% support: U12-SRC-HELP-14
% mdparagraph:P446:begin
\textbf{14. slōgon ealle and cwealdon ðæt hīe gemētton.} Cwealdon means killed. Slōgon and cwealdon reinforce the attack. Ealle is all, naturally people in context; the following neuter ðæt is a headless/\allowbreak{}collective all that they encountered. A translation all they encountered need not pretend that ðæt is an overt masculine personal relative. Cook’s Latin note reads cædunt.
% mdparagraph:P446:end

% teaching: NEW-U12-GRAM-NEW-U12-SRC-15
% teaching: NEW-U12-REUSED-U12-SRC-HELP-15
% support: U12-SRC-HELP-15
% mdparagraph:P447:begin
\textbf{15. swā-swā rīpe ierðe fortrǣdon and fornōmon.} Swā-swā means just as/\allowbreak{}like. Rīpe ierðe means ripe crops, a collective feminine accusative phrase. Fortrǣdon, trampled down; fornōmon, destroyed/\allowbreak{}laid waste. The Britons are treated as a crop that is trampled; the object continues from the preceding clauses.
% mdparagraph:P447:end

% teaching: NEW-U12-GRAM-NEW-U12-SRC-16
% support: U12-SRC-HELP-16
% mdparagraph:P448:begin
\textbf{16. and hīe mid ealle foriermdon.} Foriermdon means ruined/\allowbreak{}reduced to misery. Hīe is the Britons as object. Mid ealle is the fixed phrase completely/\allowbreak{}utterly; ealle is instrumental here. The complete phrase is supplied.
% mdparagraph:P448:end

\reserve{7\baselineskip}
% mdparagraph:P449:begin
\subsection*{S01}
% mdparagraph:P449:end

\par\begin{minipage}{\linewidth}\raggedright
% item: NEW-U12-S01
% mdparagraph:P450:begin
Translate the complete Picts and Scots extract, using the printed local helps and original notes. Preserve the renewed attack after the Roman withdrawal.
% mdparagraph:P450:end

\end{minipage}\par\addvspace{4pt}

\reserve{7\baselineskip}
% mdparagraph:P451:begin
\subsection*{S02}
% mdparagraph:P451:end

\par\begin{minipage}{\linewidth}\raggedright
% item: NEW-U12-S02
% mdparagraph:P452:begin
Put the narrated stages in order: renewed attack; appeal to Rome; first raids; Roman withdrawal; Roman relief. Explain the cause of the renewed attack as the narrator presents it.
% mdparagraph:P452:end

\end{minipage}\par\addvspace{4pt}

\reserve{7\baselineskip}
% mdparagraph:P453:begin
\subsection*{S03}
% mdparagraph:P453:end

\par\begin{minipage}{\linewidth}\raggedright
% item: NEW-U12-S03
% mdparagraph:P454:begin
In Ðā ðæt ðā ongēaton ðā ǣrran gewinnan, identify the three occurrences of ðā and the role of the first ðæt. Use the printed source help; distinguish grammatical form from clause function.
% mdparagraph:P454:end

\end{minipage}\par\addvspace{4pt}

\reserve{7\baselineskip}
% mdparagraph:P455:begin
\subsection*{S04}
% mdparagraph:P455:end

\par\begin{minipage}{\linewidth}\raggedright
% item: NEW-U12-S04
% mdparagraph:P456:begin
Who sends messengers to Rome, who sends the army, and who fights the enemies? Explain how the changing subjects are recovered from context, using the original notes where indicated.
% mdparagraph:P456:end

\end{minipage}\par\addvspace{4pt}

\reserve{7\baselineskip}
% mdparagraph:P457:begin
\subsection*{S05}
% mdparagraph:P457:end

\par\begin{minipage}{\linewidth}\raggedright
% item: NEW-U12-S05
% mdparagraph:P458:begin
Parse wæs onweg gewiten. State which event precedes which other event, and explain why this is not an agentless passive.
% mdparagraph:P458:end

\end{minipage}\par\addvspace{4pt}

\reserve{7\baselineskip}
% mdparagraph:P459:begin
\subsection*{S06}
% mdparagraph:P459:end

\par\begin{minipage}{\linewidth}\raggedright
% item: NEW-U12-S06
% mdparagraph:P460:begin
Cook’s note on Ðā onsęndon hīe him micelne hęre points out an active/\allowbreak{}passive difference from the Latin source. Identify the Old English subject, recipient and object. Give an English passive rendering that preserves their roles; no Latin or new Old English composition is required.
% mdparagraph:P460:end

\end{minipage}\par\addvspace{4pt}

\reserve{7\baselineskip}
% mdparagraph:P461:begin
\subsection*{S07}
% mdparagraph:P461:end

\par\begin{minipage}{\linewidth}\raggedright
% item: NEW-U12-S07
% mdparagraph:P462:begin
In lǣrdon ðæt hīe fæsten worhten him tō gebeorge, what action is recommended? Is that action earlier than, simultaneous with, or prospective from the instruction? Does this clause alone state that the work was completed?
% mdparagraph:P462:end

\end{minipage}\par\addvspace{4pt}

\reserve{7\baselineskip}
% mdparagraph:P463:begin
\subsection*{S08}
% mdparagraph:P463:end

\par\begin{minipage}{\linewidth}\raggedright
% item: NEW-U12-S08
% mdparagraph:P464:begin
Explain the comparison with rīpe ierðe in the final attack. Who are the victims, and how does the comparison contribute to the account? Use the lexical help rather than assuming a new ordinary crop meaning for the people.
% mdparagraph:P464:end

\end{minipage}\par\addvspace{4pt}

\reserve{8\baselineskip}
% mdparagraph:P465:begin
\section*{13. Cook’s original notes}
% mdparagraph:P465:end

\begingroup\fontsize{10}{12.8}\selectfont
\setlength{\parskip}{1.5pt}
% mdparagraph:P466:begin
The following notes reproduce the apparatus attached to this extract. Their English commentary and questions are Cook’s editorial material, not additional exercises. In particular, p. 145 note 2 preserves nineteenth-century ethnological speculation and terminology as historical apparatus; it is not a present-day account taught by this course. The fragment at the top of p. 145 continuing p. 144 note 10 is outside the extract and is not included. Notes 16–19 on p. 146 likewise belong to later text.
% mdparagraph:P466:end

% mdparagraph:P467:begin
The Old English above remains uninterrupted. This anchor list locates the original note cues; where one note is attached twice, both anchors are given.
% mdparagraph:P467:end

% mdtable:T71:begin
\par\begingroup
\fontsize{10.5}{13.2}\selectfont
\setlength{\LTleft}{0pt}\setlength{\LTright}{\fill}
\begin{longtable}{@{}>{\raggedright\arraybackslash}p{49.35mm}>{\raggedright\arraybackslash}p{71.84mm}@{}}
\toprule
\textbf{Printed page /\allowbreak{} note} & \textbf{Attached word or phrase} \\
\midrule
\endfirsthead
\toprule
\textbf{Printed page /\allowbreak{} note} & \textbf{Attached word or phrase} \\
\midrule
\endhead
\midrule
\multicolumn{2}{r}{\footnotesize Continued on the following page.}\\
\endfoot
\bottomrule
\endlastfoot
145.1 & Ðā \\
145.2 & Pihtas \\
145.3 & ǣrendwrecan \\
145.4 & gewritum \\
145.5 & fultumes \\
145.6 & underðēodnesse \\
146.1 & fīend \\
146.2 & Ðā \\
146.3 & onsęndon \\
146.4 & hīe \\
146.5 & hiera \\
146.6 & āflīemdon \\
146.7 & fæsten \\
146.8 & him \\
146.9 & sige \\
146.10 & fōron \\
146.11 & Ðā \\
146.12 & gewinnan \\
146.13 & slōgon \\
146.13 & cwealdon \\
146.14 & ierðe \\
146.15 & ealle \\
\end{longtable}
\par\endgroup

% mdtable:T71:end
% source-note: U12-SRC-P145-N01
\reserve{5\baselineskip}
% mdparagraph:P468:begin
\textbf{Printed p. 145, note 1.}
% mdparagraph:P468:end

% mdparagraph:P469:begin
“[This account] is pieced together as an abridgment of Gildas, xi.–xvi.; but the turgidity of his style is chastened, and his faulty grammar in several places corrected” (Moberly, pp. 26, 27).
% mdparagraph:P469:end

% source-note: U12-SRC-P145-N02
\reserve{5\baselineskip}
% mdparagraph:P470:begin
\textbf{Printed p. 145, note 2.}
% mdparagraph:P470:end

% mdparagraph:P471:begin
On the Picts the last edition of Chambers’ Encyclopædia remarks: “Four hypotheses have been formed in regard to the language and origin of the Picts. The first, started by Pinkerton and put by Sir Walter Scott into the mouth of the ‘Antiquary,’ is that they were Teutons, speaking a Gothic dialect; the second, maintained by Dr. Skene, is that they were Gaelic-speaking Celts; . . . the third, due to Professor Rhys, is that the Picts were non-Aryans, whose language was overlaid by loans from Welsh and Irish; and the fourth, held by two of the most eminent Celtic scholars of the day, Professor Windisch and Dr. Whitley Stokes, is that they were Celts, but more nearly allied to the Cymry than to the Gael. . . . The conclusion to which we come is that the Picts, whatever traces they show of a non-Aryan racial element, . . . spoke a Celtic language belonging to a branch of Celtic allied to the Cymric, . . . and that this dialect of the Gallo-Cymric stock was a wave of Celtic speech from the continent previous to the Gaulish which held England when Cæsar entered Britain.”
% mdparagraph:P471:end

% source-note: U12-SRC-P145-N03
\reserve{5\baselineskip}
% mdparagraph:P472:begin
\textbf{Printed p. 145, note 3.}
% mdparagraph:P472:end

% mdparagraph:P473:begin
Lat. legatos.
% mdparagraph:P473:end

% source-note: U12-SRC-P145-N04
\reserve{5\baselineskip}
% mdparagraph:P474:begin
\textbf{Printed p. 145, note 4.}
% mdparagraph:P474:end

% mdparagraph:P475:begin
Lat. epistolis.
% mdparagraph:P475:end

% source-note: U12-SRC-P145-N05
\reserve{5\baselineskip}
% mdparagraph:P476:begin
\textbf{Printed p. 145, note 5.}
% mdparagraph:P476:end

% mdparagraph:P477:begin
See 156. b.
% mdparagraph:P477:end

% source-note: U12-SRC-P145-N06
\reserve{5\baselineskip}
% mdparagraph:P478:begin
\textbf{Printed p. 145, note 6.}
% mdparagraph:P478:end

% mdparagraph:P479:begin
This pair of phrases renders subjectionem continuam. What parallel to the use of such synonymous terms may be found in the English Prayer-Book? How is it to be accounted for?
% mdparagraph:P479:end

% source-note: U12-SRC-P146-N01
\reserve{5\baselineskip}
% mdparagraph:P480:begin
\textbf{Printed p. 146, note 1.}
% mdparagraph:P480:end

% mdparagraph:P481:begin
See 46.
% mdparagraph:P481:end

% source-note: U12-SRC-P146-N02
\reserve{5\baselineskip}
% mdparagraph:P482:begin
\textbf{Printed p. 146, note 2.}
% mdparagraph:P482:end

% mdparagraph:P483:begin
This may be anywhere between A.D. 388 and 420. See Moberly, p. 27.
% mdparagraph:P483:end

% source-note: U12-SRC-P146-N03
\reserve{5\baselineskip}
% mdparagraph:P484:begin
\textbf{Printed p. 146, note 3.}
% mdparagraph:P484:end

% mdparagraph:P485:begin
This clause translates Quibus mox legio destinatur armata. Note the use of the active for the passive, which also appears in other sentences of the context.
% mdparagraph:P485:end

% source-note: U12-SRC-P146-N04
\reserve{5\baselineskip}
% mdparagraph:P486:begin
\textbf{Printed p. 146, note 4.}
% mdparagraph:P486:end

% mdparagraph:P487:begin
The legionaries, apparently.
% mdparagraph:P487:end

% source-note: U12-SRC-P146-N05
\reserve{5\baselineskip}
% mdparagraph:P488:begin
\textbf{Printed p. 146, note 5.}
% mdparagraph:P488:end

% mdparagraph:P489:begin
Of the Britons; Lat. sociorum finibus.
% mdparagraph:P489:end

% source-note: U12-SRC-P146-N06
\reserve{5\baselineskip}
% mdparagraph:P490:begin
\textbf{Printed p. 146, note 6.}
% mdparagraph:P490:end

% mdparagraph:P491:begin
Lat. expulit. See above, p. 145, n. 6.
% mdparagraph:P491:end

% source-note: U12-SRC-P146-N07
\reserve{5\baselineskip}
% mdparagraph:P492:begin
\textbf{Printed p. 146, note 7.}
% mdparagraph:P492:end

% mdparagraph:P493:begin
Lat. murum.
% mdparagraph:P493:end

% source-note: U12-SRC-P146-N08
\reserve{5\baselineskip}
% mdparagraph:P494:begin
\textbf{Printed p. 146, note 8.}
% mdparagraph:P494:end

% mdparagraph:P495:begin
See 184. a.
% mdparagraph:P495:end

% source-note: U12-SRC-P146-N09
\reserve{5\baselineskip}
% mdparagraph:P496:begin
\textbf{Printed p. 146, note 9.}
% mdparagraph:P496:end

% mdparagraph:P497:begin
Lat. triumpho. See 175.
% mdparagraph:P497:end

% source-note: U12-SRC-P146-N10
\reserve{5\baselineskip}
% mdparagraph:P498:begin
\textbf{Printed p. 146, note 10.}
% mdparagraph:P498:end

% mdparagraph:P499:begin
A passage of the Latin is here omitted in the translation, describing the construction of the (earthen) wall, between the Friths of Forth and Clyde.
% mdparagraph:P499:end

% source-note: U12-SRC-P146-N11
\reserve{5\baselineskip}
% mdparagraph:P500:begin
\textbf{Printed p. 146, note 11.}
% mdparagraph:P500:end

% mdparagraph:P501:begin
These three ðā’s respectively = when, then, and the.
% mdparagraph:P501:end

% source-note: U12-SRC-P146-N12
\reserve{5\baselineskip}
% mdparagraph:P502:begin
\textbf{Printed p. 146, note 12.}
% mdparagraph:P502:end

% mdparagraph:P503:begin
Lat. inimici.
% mdparagraph:P503:end

% source-note: U12-SRC-P146-N13
\reserve{5\baselineskip}
% mdparagraph:P504:begin
\textbf{Printed p. 146, note 13.}
% mdparagraph:P504:end

% mdparagraph:P505:begin
Lat. cædunt. See above, p. 145, n. 6.
% mdparagraph:P505:end

% source-note: U12-SRC-P146-N14
\reserve{5\baselineskip}
% mdparagraph:P506:begin
\textbf{Printed p. 146, note 14.}
% mdparagraph:P506:end

% mdparagraph:P507:begin
Lat. segetem.
% mdparagraph:P507:end

% source-note: U12-SRC-P146-N15
\reserve{5\baselineskip}
% mdparagraph:P508:begin
\textbf{Printed p. 146, note 15.}
% mdparagraph:P508:end

% mdparagraph:P509:begin
Mid ealle = completely. See 175.
% mdparagraph:P509:end

% mdparagraph:P510:begin
The transcription preserves Cook’s letters, written lengths and lexical spellings. Line wrapping and line-end division have been joined; original note cues are represented by the separate anchor list. \textbf{Rome}, \textbf{onweg}, \textbf{hęre}, \textbf{ādrifon} and \textbf{Rōmanisca} retain their source spellings. Page 146 note 13 reads \textbf{Lat. cædunt}. The supplied modern explanations in §12 are separate from the original notes.
% mdparagraph:P510:end

\reserve{7\baselineskip}
% mdparagraph:P511:begin
\section*{14. Sources for the lesson}
% mdparagraph:P511:end

% mdparagraph:P512:begin
The exercise prompts are from Albert S. Cook, \emph{Exercises in Old English} (1895), Group VI, Exercises 19–21, printed pp. 9–11; the registered scan uses PDF pp. 17–19. The reading and grammar are from Cook, \emph{A First Book in Old English}, second edition (1895). In that registered scan a printed page corresponds to the one-based PDF page sixteen higher. The reading locators and exclusions are given in §§12–13.
% mdparagraph:P512:end

% mdparagraph:P513:begin
For noun and adjective morphology, compare Cook §§43–60, printed pp. 28–41, and the vocabulary entries for the individual words. Fōt and fēond are treated at §§46–46.3; Wright, \emph{Old English Grammar} (1914), §§416–417, printed p. 202 /\allowbreak{} PDF p. 222, independently supports the fēond class and genitive plural fēonda. Cook §44.2 gives the hęre alternatives; §57.3 gives the optional syncope of ēadig. His vocabulary explicitly distinguishes the noun and adjective īsern.
% mdparagraph:P513:end

% mdparagraph:P514:begin
For the defined verb systems, commands and nonfinite bases, compare Cook §§93–95 and the individual strong and weak classes; §127 for āgan and §139 for willan and its negative forms. The selected dental combinations follow Cook §§33–35 and Wright §476, printed p. 248. Cȳst/\allowbreak{}cȳð and the attested cȳdde alternative are checked with Bosworth–Toller, s.v. cȳðan.
% mdparagraph:P514:end

% mdparagraph:P515:begin
The treatment of time and reported viewpoint uses Cook §§186–188 and §§192–196, with Fulk (2020), §§45–46 and §§48–52. Bosworth–Toller, \emph{An Anglo-Saxon Dictionary}, main entry habban IV–V, printed p. 497a, and the Supplement, habban B, printed p. 493a, provide evidence for both uninflected and object-agreeing participles. The chapter teaches only the selected constructions described above.
% mdparagraph:P515:end

% mdparagraph:P516:begin
Prepositional and case uses draw on Cook §172 and §§160–161, 164, with the individual lexical evidence. Bosworth–Toller’s ēhtan entry includes the same pursuing-in-this-city expression; its gemang entry supports the assigned sheep-and-wolves comparison and accusative motion phrase. Its entries mā and swā support the quantitative and comparative distinctions. The dictionary’s mild-heort entry, printed p. 687b, supports merciful in the learning saying; stán-weall, printed p. 912a, supports stone wall. In teaching lemmas, dictionary division hyphens are joined and the evidenced acute quantity of stán-weall is written with the course’s macron, stānweall. No unmarked vowel length is supplied by guesswork.
% mdparagraph:P516:end

\endgroup
\end{document}
